\chapter{Smooth Monge--Amp\`ere equations}
\label{chap:smooth-equations}

Let $X$ be a connected compact complex manifold of dimension $n\geq1$. Fix a Hermitian form $\omega$ and a smooth positive volume form $\mathrm{d}V$ on $X$. Given $F\in C^\infty(X,\mathbb R)$, we seek a smooth real potential $u$ and a scalar $b\in\mathbb R$ satisfying
\begin{equation}\label{eq:smooth-ma}
 \omega_u^n=\mathrm{e}^{F+b}\mathrm{d}V,\qquad \omega_u>0.
\end{equation}
Here $\omega_u=\omega+\ddc u$, and adding a constant to $u$ does not change the equation. For a closed background, integration determines $b$ from the prescribed density. For a Hermitian background, the total mass can depend on $u$, so the potential and the scalar must be found together. The adjoint probability supplies their linearized compatibility condition. Weighted integration and a trace estimate give the compactness needed to continue solutions from the background density; the same adjoint probability then governs the variation of $b$.

For $n\geq2$, prescribing the volume of a Gauduchon metric requires a different potential variation: we deform its $(n-1)$st power within an Aeppli class. The resulting first-order torsion terms have a specific matrix structure. Their coefficients make the derivative estimates possible, and rescaling reduces the gradient bound to a Liouville theorem on $\mathbb C^n$.

\section{Scalar bounds and uniqueness}

Extrema of the potential bound the normalization scalar in terms of the prescribed data. Applied to the difference of two solutions, the same argument identifies their scalars; ellipticity then determines the potential up to a constant. The scalar bounds also make the continuity path uniform.

\begin{theorem}[Tosatti--Weinkove]\label{thm:tw-smooth}\index{Monge--Ampere equation@Monge--Amp\`ere equation!smooth data}
For every smooth real $F$, \eqref{eq:smooth-ma} has a unique pair $(u,b)\in C^\infty(X,\mathbb R)\times\mathbb R$ with $\sup_Xu=0$.
\end{theorem}

In dimension one, integration determines
$\mathrm{e}^b=\int_X\omega/\int_X\mathrm{e}^F \mathrm{d}V$, and the remaining equation for $u$ is the usual Poisson equation. Thus the statement holds in that dimension as well.

\begin{proposition}\label{prop:b-extrema}\index{Normalization scalar!extremum bounds}
Every solution of \eqref{eq:smooth-ma} satisfies
\[
 \inf_X\!\left(\log\frac{\omega^n}{\mathrm{d}V}-F\right)
 \leq b\leq
 \sup_X\!\left(\log\frac{\omega^n}{\mathrm{d}V}-F\right).
\]
\end{proposition}
\begin{proof}
At a maximum point $x_+$ of $u$, the Hermitian form $\ddc u$ is nonpositive. Hence $\omega_u^n(x_+)\leq\omega^n(x_+)$, and
$b\leq\log(\omega^n/\mathrm{d}V)(x_+)-F(x_+)$.
At a minimum point the reverse argument gives the lower bound.
\end{proof}

When $\mathrm{d}\omega=0$, integration gives the familiar explicit formula
\begin{equation}\label{eq:closed-b}
 b=\log\frac{\int_X\omega^n}{\int_X\mathrm{e}^F \mathrm{d}V}.
\end{equation}
The same formula holds whenever the mass-invariance criterion of \cref{thm:mass-invariance-characterization} is satisfied. In general, integrating the equation gives only
\[
 \mathrm{e}^b=\frac{\Mass_\omega(u)}{\int_X\mathrm{e}^F \mathrm{d}V},
\]
whose numerator depends on the unknown potential.

\begin{proposition}\label{prop:smooth-uniqueness}\index{Uniqueness!smooth positive equation}
Suppose $(u,b)$ and $(v,c)$ are smooth solutions of \eqref{eq:smooth-ma} with the same $F$ and $\mathrm{d}V$. Then $b=c$ and $u-v$ is constant.
\end{proposition}
\begin{proof}
At a maximum of $u-v$, one has $\omega_u\leq\omega_v$, so $\mathrm{e}^b\leq \mathrm{e}^c$. At a minimum one obtains $\mathrm{e}^b\geq \mathrm{e}^c$. Therefore $b=c$.

Put $h=u-v$ and $\eta_t=\omega_v+t\ddc h$ for $0\leq t\leq1$. This is a family of strictly positive forms. Subtracting the logarithms of the two equations and integrating the derivative along the segment gives
\[
 0=\int_0^1 n\,\frac{\ddc h\wedge\eta_t^{n-1}}{\eta_t^n}\,\mathrm{d}t.
\]
The right side is a linear elliptic operator applied to $h$. Its coefficients are smooth and uniformly positive on the compact manifold, and it has no zeroth-order term. The strong maximum principle implies that $h$ is constant.
\end{proof}

\section{The potential bound from weighted integration}

An upper bound for the density controls the normalized potential. Integrating against $\mathrm{e}^{-pu}$ produces a gradient term and mixed weighted masses; a finite recurrence absorbs the latter together with the torsion error. Sobolev iteration then gives a set of uniformly positive volume near the minimum, and normalized $L^1$ compactness bounds that minimum. Norms below use the fixed metric $\omega$.

\Needspace{10\baselineskip}
\begin{lemma}\label{lem:pointwise-torsion-estimate}
There is a constant $C$ depending on $\omega$ such that, for every positive $(1,1)$-form $\eta$, every smooth real $f$, every $\delta>0$, and $0\leq k\leq n-2$,
\[
 \begin{split}
 |\dc f\wedge \mathrm{d}\omega\wedge\eta^k\wedge\omega^{n-k-2}|
 \leq{}&\frac C\delta\,\mathrm{d}f\wedge\dc f\wedge\eta^k\wedge\omega^{n-k-1}\\
 &+C\delta\,\eta^k\wedge\omega^{n-k}.
 \end{split}
\]
The absolute value means that of the coefficient relative to the positive volume form $\omega^n$.
\end{lemma}
\begin{proof}
At the point in question choose an $\omega$-unitary frame in which $\eta$ is diagonal, with positive eigenvalues $\lambda_1,\ldots,\lambda_n$. Expand $\eta^k$ as a sum over $k$-element subsets $I$, with coefficient a positive dimensional multiple of $\lambda_I=\prod_{i\in I}\lambda_i$. A surviving coefficient from $\dc f\wedge \mathrm{d}\omega$ has the form of a bounded coefficient of $\mathrm{d}\omega$ times one first derivative $f_a$, where $a\notin I$. Young's inequality bounds its absolute value by
\[
 C\lambda_I|f_a|\leq (C/\delta)\lambda_I|f_a|^2
                              +C\delta\lambda_I.
\]
The sum of the first terms is bounded by a fixed multiple of the first positive form in the assertion: its diagonal expansion contains precisely the terms $\lambda_I|f_a|^2$ with $a\notin I$. The second positive form contains all the $\lambda_I$. Finitely many indices and the uniform bound for $\mathrm{d}\omega$ give a uniform $C$.
\end{proof}

\begin{lemma}\label{lem:cherrier-gradient}
Suppose $u$ is smooth, $\eta=\omega+\ddc u>0$, and
$\eta^n\leq\Lambda\omega^n$ for a fixed $\Lambda>0$. There are constants $p_0\geq1$ and $C$ depending only on $\omega$ and $\Lambda$ such that
\[
 \int_X\mathrm{d}(\mathrm{e}^{-pu/2})\wedge\dc(\mathrm{e}^{-pu/2})\wedge\omega^{n-1}
 \leq Cp\int_X\mathrm{e}^{-pu}\omega^n\quad(p\geq p_0).
\]
\end{lemma}
\begin{proof}
For a fixed real $p\geq1$, define $M_k$ for $0\leq k\leq n$ and $J_k$ for $0\leq k\leq n-1$ by
\[
 M_k=\int_X\mathrm{e}^{-pu}\eta^k\wedge\omega^{n-k},\qquad
 J_k=\int_X\mathrm{e}^{-pu}\mathrm{d}u\wedge\dc u\wedge\eta^k\wedge\omega^{n-k-1},
 \quad J=\sum_{k=0}^{n-1}J_k.
\]
All these numbers are nonnegative. Let
$T=\sum_{k=0}^{n-1}\eta^k\wedge\omega^{n-k-1}$. Since $\mathrm{d}\eta=\mathrm{d}\omega$, differentiation gives
\[
 \mathrm{d}T=n\sum_{k=0}^{n-2}\eta^k\wedge\omega^{n-k-2}\wedge \mathrm{d}\omega.
\]
The sum is empty when $n=1$. Stokes' theorem gives the exact identity
\[
 \begin{split}
 \int_X\mathrm{e}^{-pu}(\eta^n-\omega^n)
 &=\int_X\mathrm{e}^{-pu}\ddc u\wedge T\\
 &=pJ+\int_X\mathrm{e}^{-pu}\dc u\wedge \mathrm{d}T.
 \end{split}
\]
The density assumption bounds the left side above by $(\Lambda+1)M_0$. Apply \cref{lem:pointwise-torsion-estimate} to the last integral, and choose $\delta=K/p$ for a sufficiently large fixed $K$. The contribution proportional to $J/\delta$ is then at most $pJ/2$. After absorption we obtain
\begin{equation}\label{eq:weighted-first-bound}
 pJ\leq CM_0+\frac Cp\sum_{k=0}^{n-2}M_k.
\end{equation}
This estimate alone is not enough, because the mixed weighted masses on the right have not been bounded.

For $1\leq k\leq n-2$, a second integration by parts gives
\begin{equation}\label{eq:weighted-recurrence}
 M_k-M_{k-1}=pJ_{k-1}+R_k,
 \qquad
 R_k=\int_X\mathrm{e}^{-pu}\dc u\wedge
                \mathrm{d}(\eta^{k-1}\wedge\omega^{n-k}).
\end{equation}
The derivative in $R_k$ has two terms, with coefficients $k-1$ and $n-k$. Apply the pointwise estimate with $\delta=1/p$ to each of them. Omitting nonexistent negative indices yields
\[
 |R_k|\leq Cp(J_{k-2}+J_{k-1})
                   +\frac Cp(M_{k-2}+M_{k-1}),
 \qquad M_{-1}=J_{-1}=0.
\]
Consequently, for $p\geq1$,
\[
 M_k\leq C(M_{k-1}+M_{k-2}+pJ).
\]
Induction over the finitely many values of $k$ proves
\[
 \sum_{k=0}^{n-2}M_k\leq C(M_0+pJ).
\]
For $n=2$ this is immediate, and for $n=1$ the sum is empty. Inserting this bound in \eqref{eq:weighted-first-bound} gives
$pJ\leq CM_0+C'J$. Taking $p\geq2C'$ absorbs the last term and proves $pJ\leq CM_0$.

Finally,
\[
 \int_X\mathrm{d}(\mathrm{e}^{-pu/2})\wedge\dc(\mathrm{e}^{-pu/2})\wedge\omega^{n-1}
 =\frac{p^2}{4}J_0\leq\frac{p^2}{4}J\leq CpM_0.
\]
\end{proof}

\Needspace{10\baselineskip}
\begin{theorem}\label{thm:smooth-czero}\index{A priori estimate!potential bound}
For each $\Lambda>0$ there is $C_\Lambda$ such that every smooth $u$ with
\[
 \omega+\ddc u>0,\qquad \sup_Xu=0,\qquad
 (\omega+\ddc u)^n\leq\Lambda\omega^n
\]
satisfies $\|u\|_\infty\leq C_\Lambda$.
\end{theorem}
\Needspace{8\baselineskip}
\begin{proof}
Let $\mathrm{d}\mu=\omega^n/\int_X\omega^n$. The usual Sobolev inequality for the fixed Riemannian metric of $\omega$ gives, for some $\beta>1$ and every smooth real test function $w$,
\[
 \left(\int_X |w|^{2\beta}\,\mathrm{d}\mu\right)^{1/\beta}
 \leq C\left(\int_X\mathrm{d}w\wedge\dc w\wedge\omega^{n-1}
                         +\int_X|w|^2\,\mathrm{d}\mu\right).
\]
We may take $\beta=n/(n-1)$ if $n\geq2$, and for example $\beta=2$ if $n=1$. Choose $p_0$ as in \cref{lem:cherrier-gradient}. Apply the Sobolev inequality to $w=\mathrm{e}^{-pu/2}$ and use that lemma. With $Z=\mathrm{e}^{-u}$ we get
\[
 \|Z\|_{L^{p\beta}(\mu)}\leq(Cp)^{1/p}\|Z\|_{L^p(\mu)}
 \quad(p\geq p_0).
\]
Iterate at $p=p_0\beta^j$. The product of the constants is finite because
$\sum_j(1+\log(p_0\beta^j))/(p_0\beta^j)<\infty$. Hence, for a uniform $C_1\geq1$ and $m=\inf_Xu$,
\begin{equation}\label{eq:near-minimum-integral}
 \mathrm{e}^{-p_0m}\leq C_1\int_X\mathrm{e}^{-p_0u}\,\mathrm{d}\mu.
\end{equation}
This still involves the unknown potential on the right. Divide by $\mathrm{e}^{-p_0m}$ and set $A=p_0^{-1}\log(2C_1)$. The integrand $\mathrm{e}^{-p_0(u-m)}$ is at most one everywhere and at most $(2C_1)^{-1}$ outside $S=\{u\leq m+A\}$. Therefore \eqref{eq:near-minimum-integral} implies
\[
 \mu(S)\geq\frac1{2C_1}=:\delta>0.
\]
By \cref{thm:normalized-compactness}, $\int_X(-u)\,\mathrm{d}\mu\leq C_0$. On $S$ we have $-u\geq-m-A$, so
\[
 C_0\geq\int_X(-u)\,\mathrm{d}\mu\geq\delta(-m-A)
\]
when $-m\geq A$; otherwise the desired bound is immediate. Thus $-m\leq A+C_0/\delta$, proving the theorem.
\end{proof}

\section{The linearized equation and smooth dependence}

The linearized operator $L_u$ of \eqref{eq:linearized-ma} annihilates constants, and its range is the space of functions with zero mean under the adjoint probability $\mu_u$ of \cref{prop:adjoint-gauduchon}. Including the scalar as an unknown removes this obstruction. In dimension one, $\mu_u=\omega_u/\int_X\omega_u$; the compatibility condition follows from Stokes' theorem.

\Needspace{7\baselineskip}
The normalization $\sup_Xu=0$ is convenient for compactness, but need not vary differentiably with the data. For differentiating solutions, fix instead the linear normalization
\[
 \int_Xu\,\mathrm{d}V=0.
\]
Changing between these normalizations only adds a constant to $u$ and leaves $b$ unchanged.

\begin{proposition}\label{prop:smooth-dependence}\index{Monge--Ampere equation@Monge--Amp\`ere equation!smooth dependence}
With this linear normalization, the solution $(u,b)$ of \eqref{eq:smooth-ma} depends smoothly on $F$ in H\"older spaces, and hence along every smooth finite-dimensional family of smooth data.
\end{proposition}
\begin{proof}
Fix an integer $k\geq0$ and a real exponent $0<\alpha<1$. Consider the map
\[
 (u,b,F)\longmapsto
 \log\frac{(\omega+\ddc u)^n}{\mathrm{d}V}-F-b
\]
on the open set of potentials in $C^{k+2,\alpha}$ with $\omega_u>0$, restricted to $\int_Xu\,\mathrm{d}V=0$. Its derivative in $(u,b)$ is
\[
 (v,a)\longmapsto L_uv-a.
\]
Given $h\in C^{k,\alpha}$, integrating $L_uv-a=h$ against $\mu_u$ forces
$a=-\int_Xh\,\mathrm{d}\mu_u$. For this value, the equation for $v$ satisfies the compatibility condition of \cref{prop:adjoint-gauduchon}, and has a unique solution with zero $\mathrm{d}V$-mean. Elliptic estimates make this inverse bounded. The Banach-space implicit function theorem gives a local solution map around the given positive solution; this step does not require the global existence theorem. Uniqueness identifies the local maps wherever they overlap, and elliptic regularity allows arbitrary $k$ for smooth data.
\end{proof}

\section{Second derivatives and ellipticity}

We now supply the estimate that turns a bound for the potential into uniform ellipticity. The torsion calculation goes back to Cherrier. We use its completed-square form and the maximum argument of Tosatti--Weinkove \cite[Sections 3 and 9]{TW15}.

For this calculation write
\[
 \omega=\frac{\mathrm i}{2\pi}g_{i\bar j}\,\mathrm{d}z^i\wedge \mathrm{d}\bar z^j,
 \qquad \omega_u=\frac{\mathrm i}{2\pi}h_{i\bar j}\,\mathrm{d}z^i\wedge \mathrm{d}\bar z^j,
 \qquad h_{i\bar j}=g_{i\bar j}+u_{i\bar j}.
\]
Thus $\Delta_gv=g^{i\bar j}v_{i\bar j}$ and $\Delta_hv=h^{i\bar j}v_{i\bar j}$ are exactly the trace operators for our $\ddc$ convention. In particular, $\Delta_hu=n-\operatorname{tr}_hg$. Repeated indices are summed. Let $\nabla$ be the Chern connection of the fixed metric $g$, with
\[
 \Gamma_{ik}^p=g^{p\bar q}\partial_i g_{k\bar q},\qquad
 T_{ik}^p=\Gamma_{ik}^p-\Gamma_{ki}^p,\qquad
 \tau_k=\sum_iT_{ik}^i.
\]
All bounds for this connection and its curvature depend only on $g$.

\begin{lemma}\label{lem:hermitian-trace-inequality}
Suppose $h=g+(u_{i\bar j})>0$, and put
\[
 S=\operatorname{tr}_gh,\qquad B=\operatorname{tr}_hg,
 \qquad G=\log\frac{\det h}{\det g}.
\]
There is $C_g$ such that, at every point where $S\geq1$,
\begin{equation}\label{eq:hermitian-trace-inequality}
 \Delta_h\log S\geq
 \frac2{S^2}\operatorname{Re}\bigl(h^{k\bar q}\tau_k S_{\bar q}\bigr)
 +\frac{\Delta_gG}{S}-C_gB.
\end{equation}
\end{lemma}
\begin{proof}
The identity $\mathrm{d}\omega_u=\mathrm{d}\omega$ is equivalent to
\begin{equation}\label{eq:trace-torsion-switch}
 \nabla_i h_{k\bar j}-\nabla_k h_{i\bar j}
 =T_{ik}^p(g_{p\bar j}-h_{p\bar j}),
\end{equation}
and its complex conjugate. This also follows directly by cancelling the third derivatives of $u$.

For tensors with two holomorphic and one antiholomorphic lower indices, use the positive norm
\[
 \|A\|^2=g^{i\bar\ell}h^{k\bar q}h^{p\bar j}
 A_{ik\bar j}\overline{A_{\ell q\bar p}}.
\]
Set $A_{ik\bar j}=\nabla_i h_{k\bar j}$ and
$D_{ik\bar j}=T_{ik}^p h_{p\bar j}$. We first claim
\begin{equation}\label{eq:trace-before-square}
 \Delta_hS=\|A+D\|^2+\Delta_gG+E,
 \qquad |E|\leq C_g(1+S)B.
\end{equation}
Here is the calculation and the bound for its remainder. Start with
$\Delta_hS=h^{i\bar j}g^{k\bar\ell}\nabla_i\nabla_{\bar j}h_{k\bar\ell}$.
Switch $\bar j,\bar\ell$ by the conjugate of \eqref{eq:trace-torsion-switch}, commute $\nabla_i,\nabla_{\bar\ell}$, and then switch $i,k$ by \eqref{eq:trace-torsion-switch}. The terms with a derivative of $h$ are
\[
 h^{i\bar j}g^{k\bar\ell}\nabla_{\bar\ell}\nabla_kh_{i\bar j}
 -2\operatorname{Re}\bigl(h^{i\bar j}g^{k\bar\ell}
             T_{ik}^p\nabla_{\bar\ell}h_{p\bar j}\bigr).
\]
In making the first switch there is also a term
$-h^{i\bar j}g^{k\bar\ell}\overline{T_{j\ell}^q}\nabla_i h_{k\bar q}$.
Using \eqref{eq:trace-torsion-switch} on it is what gives one of the two conjugate derivative terms just displayed, together with
\[
 h^{i\bar j}g^{k\bar\ell}T_{ik}^p\overline{T_{j\ell}^q}
                   (h_{p\bar q}-g_{p\bar q}).
\]
The part containing $h$ is $\|D\|^2$. Differentiating the logarithmic determinant gives
\[
 \Delta_gG=h^{i\bar j}g^{k\bar\ell}
                \nabla_{\bar\ell}\nabla_kh_{i\bar j}-\|A\|^2.
\]
The two derivative terms left over are $2\operatorname{Re}\langle A,D\rangle$, using the antisymmetry of $T$. These identities give \eqref{eq:trace-before-square}.

More explicitly, let
\[
 \mathcal C_{i\bar\ell k\bar j}
  =(\nabla_i\nabla_{\bar\ell}-\nabla_{\bar\ell}\nabla_i)h_{k\bar j}.
\]
The remainder is the contraction by $h^{i\bar j}g^{k\bar\ell}$ of
\[
 \begin{split}
 &(\nabla_i\overline{T_{j\ell}^q})(g_{k\bar q}-h_{k\bar q})
 +(\nabla_{\bar\ell}T_{ik}^p)(g_{p\bar j}-h_{p\bar j})\\
 &\hspace{15mm}+\mathcal C_{i\bar\ell k\bar j}
              -T_{ik}^p\overline{T_{j\ell}^q}g_{p\bar q}.
 \end{split}
\]
The commutator $\mathcal C$ is curvature acting linearly on $h$; this follows by inserting $\nabla=\partial+\Gamma$ and cancelling the derivatives of $h$. In a $g$-unitary frame diagonalizing $h$, every coefficient inside this remainder is bounded by $C_g(1+S)$. The contraction is consequently bounded by $C_g(1+S)\sum_i h^{i\bar i}$, as asserted. No derivative of $h$ remains in $E$.

It remains to control the gradient introduced by the logarithm. Define
\[
 H_{ik\bar j}=h_{i\bar j}\frac{S_k}{S},\qquad
 K=\|A+D-H\|^2\geq0.
\]
By \eqref{eq:trace-torsion-switch},
\[
 g^{i\bar\ell}(A+D)_{ik\bar\ell}=S_k+\tau_k.
\]
Also $\|H\|^2=|\partial S|_h^2/S$. Expanding the square therefore gives
\[
 K=\|A+D\|^2-\frac{|\partial S|_h^2}{S}
       -\frac2S\operatorname{Re}(h^{k\bar q}\tau_kS_{\bar q}).
\]
Combine this with \eqref{eq:trace-before-square} to obtain the exact identity
\[
 \Delta_h\log S=\frac KS+
 \frac2{S^2}\operatorname{Re}(h^{k\bar q}\tau_kS_{\bar q})
 +\frac{\Delta_gG}{S}+\frac ES.
\]
Since $S\geq1$, the remainder is at least $-2C_gB$. Enlarging the constant proves the lemma.
\end{proof}

\Needspace{12\baselineskip}
\begin{theorem}\label{thm:hermitian-second-estimate}
Fix $D,K>0$. Suppose $u$ is smooth, $\omega_u>0$, and
\[
 \operatorname{osc}_Xu\leq D,\qquad
 \omega_u^n=e^G\omega^n,\qquad
 |G|\leq K,\qquad \Delta_gG\geq-K.
\]
There is a constant $C=C(\omega,D,K)$ such that
\[
 C^{-1}\omega\leq\omega_u\leq C\omega.
\]
\end{theorem}
\begin{proof}
Write $s=u-\inf_Xu+1$, so $1\leq s\leq D+1$, and maximize
\[
 Q=\log S-Au+s^{-1}
\]
on $X$, with $A\geq1$ to be fixed. At a maximum point $x$, if $S(x)<1$ the upper bound for $S$ everywhere already follows from the oscillation bound and the definition of $Q$. We may therefore assume $S(x)\geq1$ and use \cref{lem:hermitian-trace-inequality}.

The positive eigenvalues of $h$ relative to $g$ give $SB\geq n^2$. Hence $\Delta_gG/S\geq-KB/n^2$. Moreover, $\partial Q=0$ gives
\[
 \frac{S_k}{S}=(A+s^{-2})u_k.
\]
Because $|\tau|_h^2\leq C_gB$, the absolute value of the torsion term is bounded by
\begin{equation}\label{eq:trace-torsion-absorption}
 \begin{split}
 \left|\frac2{S^2}\operatorname{Re}(h^{k\bar q}\tau_kS_{\bar q})\right|
 &\leq\frac{C_g(A+1)}S\sqrt B\,|\partial u|_h\\
 &\leq\frac{|\partial u|_h^2}{s^3}
          +\frac{C_g(A+1)^2s^3}{S^2}B.
 \end{split}
\end{equation}
If $S(x)^2\leq(A+1)^2(D+1)^3$, then $S(x)$ is bounded once $A$ is fixed, and again so is $S$ everywhere. In the remaining case the last coefficient of $B$ is at most $C_g$.

On the other hand, direct differentiation gives
\[
 \Delta_hQ=\Delta_h\log S+(A+s^{-2})(B-n)
                         +2s^{-3}|\partial u|_h^2.
\]
The trace inequality and \eqref{eq:trace-torsion-absorption} imply at $x$
\[
 0\geq\Delta_hQ\geq(A-C_g-K/n^2)B-n(A+1),
\]
after a fixed enlargement of $C_g$. Choose $A>C_g+K/n^2+1$. It follows that $B(x)$ is bounded. If $\lambda_i$ are the eigenvalues, then for $n\geq2$
\[
 S=\sum_i\lambda_i
 =e^G\sum_i\prod_{j\ne i}\lambda_j^{-1}
 \leq n e^K B^{n-1}.
\]
For $n=1$, $S=e^G$ directly. Thus $S(x)$ is bounded in all cases. For every $y\in X$, the maximum property gives
\[
 \log S(y)\leq\log S(x)+A(u(y)-u(x))+s(x)^{-1}-s(y)^{-1}
             \leq\log S(x)+AD+1.
\]
This proves $S\leq C$ on $X$. The determinant lower bound then gives
$\lambda_i\geq e^{-K}C^{-(n-1)}$, proving uniform equivalence.
\end{proof}

\begin{corollary}\label{thm:classical-higher-estimate}
Fix $\omega$. For smooth solutions of
\[
 (\omega+\ddc u)^n=e^G\omega^n,\qquad\omega+\ddc u>0,
\]
a uniform $C^0$ bound for $u$ and uniform bounds for $G$ in every $C^k$ norm imply uniform bounds for $u$ in every $C^k$ norm.
\end{corollary}
\begin{proof}
\Cref{thm:hermitian-second-estimate} gives uniform ellipticity and a bound for $\Delta_gu=S-n$. The standard linear elliptic $W^{2,p}$ estimate for this fixed operator first bounds $u$ in $C^{1,\alpha}$ for any $\alpha<1$. The local complex Evans--Krylov estimate for the concave logarithmic determinant then gives a uniform $C^{2,\alpha}$ bound. One can use its real uniformly elliptic formulation: project real symmetric Hessians to their complex Hermitian part and extend the logarithmic determinant concavely from the compact set of positive matrices supplied by the preceding theorem. The extension is uniformly elliptic on real symmetric matrices because that projection preserves positivity and its trace. The usual Evans--Krylov theorem applies to this extension. Finally, differentiation of the equation gives linear uniformly elliptic equations with H\"older coefficients; Schauder estimates and induction give all higher bounds. Only these standard local elliptic regularity theorems are used here.
\end{proof}

The same conclusion holds for equations $\omega_u^n=e^{a u+F}\omega^n$ with $a$ in a fixed bounded interval of nonnegative numbers, $u$ uniformly bounded, and $F$ uniformly bounded in every $C^k$. Indeed,
\[
 \Delta_g(a u+F)=a(S-n)+\Delta_gF\geq-an+\Delta_gF,
\]
so the second-derivative estimate applies without a prior bound for derivatives of $u$. The fixed-operator estimate first gives $C^{1,\alpha}$, after which the same local regularity and bootstrapping apply.


\section{The continuity method}

The local inverse gives openness of the solvable set. The extrema bound for the scalar, the potential estimate and the derivative estimates give compactness along a path from the background density to the prescribed one. Uniform ellipticity preserves positivity at a limiting parameter.

\begin{proof}[of \cref{thm:tw-smooth}]
The dimension-one argument was given at the start of the chapter. In higher dimension put
$G=F+\log(\mathrm{d}V/\omega^n)$ and consider
\[
 (\omega+\ddc u_t)^n=\mathrm{e}^{tG+b_t}\omega^n,
 \qquad\sup_Xu_t=0,\qquad0\leq t\leq1.
\]
At $t=0$, take $u_0=0$ and $b_0=0$. The local implicit-function argument in \cref{prop:smooth-dependence}, with the additive constant temporarily fixed by a linear mean, shows that the set of solvable parameters is open.

At every solvable parameter, \cref{prop:b-extrema} gives
\[
 -\sup_X(tG)\leq b_t\leq-\inf_X(tG),
 \qquad \mathrm{e}^{tG+b_t}\leq \mathrm{e}^{\operatorname{osc}_XG}.
\]
The constants $b_t$ are uniformly bounded, and \cref{thm:smooth-czero} bounds $u_t$ in $C^0$. Apply \cref{cor:classical-higher-estimate}. Along any sequence of solvable parameters converging in $[0,1]$, a subsequence of $(u_t,b_t)$ therefore converges smoothly to a solution; uniform equivalence of the metrics preserves strict positivity. The supremum normalization also passes to the limit. Thus the solvable set is closed. Connectedness of $[0,1]$ gives a solution at $t=1$, and \cref{prop:smooth-uniqueness} gives uniqueness.
\end{proof}

\section{The smooth exponential equation}

When the logarithmic density contains a positive multiple of the potential, extrema bound the potential itself and the linearization acquires a negative zeroth-order term. The continuity method therefore has neither an additive normalization nor an extra scalar unknown.

\begin{proposition}\label{prop:smooth-proper-comparison}\label{prop:exp-smooth-unique}
Let $\lambda>0$, let $F$ be smooth, and let $v,w$ be smooth strictly $\omega$-plurisubharmonic functions. If
\[
 \omega_v^n\geq \mathrm{e}^{\lambda v+F}\omega^n,
 \qquad \omega_w^n\leq \mathrm{e}^{\lambda w+F}\omega^n,
\]
then $v\leq w$. In particular, two smooth solutions coincide.
\end{proposition}
\begin{proof}
If $v-w$ had a positive maximum at $x$, then $\omega_v(x)\leq\omega_w(x)$. Taking determinants and applying the two assumptions would give
$\mathrm{e}^{\lambda v(x)+F(x)}\leq \mathrm{e}^{\lambda w(x)+F(x)}$, contradicting $v(x)>w(x)$. For two solutions, apply the comparison in both orders.
\end{proof}

\begin{theorem}\label{thm:smooth-proper-existence}\index{Monge--Ampere equation@Monge--Amp\`ere equation!smooth exponential equation}
For every $\lambda>0$ and every smooth real $F$, there is a unique smooth $u$ with $\omega_u>0$ such that
\[
 \omega_u^n=\mathrm{e}^{\lambda u+F}\omega^n.
\]
The solution satisfies
\[
 -\lambda^{-1}\sup_XF\leq u\leq-\lambda^{-1}\inf_XF.
\]
\end{theorem}
\begin{proof}
At a minimum of $u$, the equation gives $\lambda u+F\geq0$; at a maximum it gives $\lambda u+F\leq0$. These inequalities prove the stated bounds.

Consider the equations
\[
 \omega_{u_t}^n=\mathrm{e}^{\lambda u_t+tF}\omega^n,\qquad0\leq t\leq1.
\]
At $t=0$, $u_0=0$ is a solution. The linearization in $u$ is $L_u-\lambda$. Its kernel is zero by the maximum principle. Its Fredholm index is zero: deform its positive principal coefficients to those of a fixed Laplacian while keeping the negative zeroth-order term. The latter operator is invertible, and elliptic Fredholm index is constant under this deformation. Thus $L_u-\lambda$ is an isomorphism between the usual H\"older spaces, with bounded inverse by Schauder theory. The implicit function theorem proves openness.

The extrema bounds give $\|u_t\|_\infty\leq\lambda^{-1}\|F\|_\infty$. The semilinear consequence of \cref{thm:hermitian-second-estimate} gives uniform ellipticity and all higher bounds. A subsequence along any converging sequence of solvable parameters therefore converges smoothly to a strictly positive solution at the limit. This proves closedness, and hence existence at $t=1$. Uniqueness is \cref{prop:exp-smooth-unique}.
\end{proof}

For a fixed smooth real obstacle $f$ on $X$ and a parameter $\beta>0$, the choice $\lambda=\beta$ and $F=-\beta f$ gives
\[
 \omega_{u_\beta}^n=\mathrm{e}^{\beta(u_\beta-f)}\omega^n.
\]
These equations penalize values above the obstacle. In \cref{prop:smooth-obstacle-support} we use them to construct its envelope and prove contact support. The necessary convergence will use only bounded regularization and the comparison just proved.

\section{Variation of the normalization scalar}

Integration against the adjoint probability eliminates the derivatives of the potential from the first variation. In the second variation, the matrix logarithm leaves a nonnegative square with a minus sign. This gives concavity and identifies its equality case.

Write $b(F)$ for the scalar in \eqref{eq:smooth-ma}, and use the linear normalization $\int_Xu\,\mathrm{d}V=0$. For a smooth family $F_t$, with $t$ in an open interval $I\subset\mathbb R$, write $(u_t,b_t)$ for its solutions and put $\eta_t=\omega_{u_t}$. A dot denotes differentiation in $t$.

\begin{theorem}\label{thm:b-derivatives}\index{Normalization scalar!smooth variation}
Put $A_t=\eta_t^{-1}\ddc\dot u_t$, the Hermitian endomorphism determined by the indicated forms. The normalization constant satisfies
\begin{align}
 \dot b_t&=-\int_X\dot F_t\,\mathrm{d}\mu_{u_t},\label{eq:b-first}\\
 \ddot b_t&=-\int_X\ddot F_t\,\mathrm{d}\mu_{u_t}
 -\int_X\operatorname{tr}(A_t^2)\,\mathrm{d}\mu_{u_t},\label{eq:b-second}
\end{align}
In particular, for smooth real functions $F_0,G$ on $X$ and the affine family $F_t=F_0+tG$, the function $b_t$ is concave. Its second derivative vanishes at a parameter value if and only if $G$ is constant on $X$.
\end{theorem}
\begin{proof}
Differentiating the logarithm of \eqref{eq:smooth-ma} gives
\[
 L_{u_t}\dot u_t=\dot F_t+\dot b_t.
\]
Integration against $\mu_{u_t}$ proves \eqref{eq:b-first}.

For positive matrices $H_t$, direct differentiation of the inverse gives
\[
 \frac{\mathrm{d}^2}{\mathrm{d}t^2}\log\det H_t
 =\operatorname{tr}(H_t^{-1}\ddot H_t)
  -\operatorname{tr}\bigl((H_t^{-1}\dot H_t)^2\bigr).
\]
Apply this pointwise to $\eta_t$. The second derivative of the equation is
\[
 L_{u_t}\ddot u_t-\operatorname{tr}(A_t^2)
 =\ddot F_t+\ddot b_t.
\]
At this fixed parameter, integrate against $\mu_{u_t}$ once more. This proves \eqref{eq:b-second}; there is no need to differentiate the measure.

The eigenvalues of the Hermitian endomorphism $A_t$ are real, so $\operatorname{tr}(A_t^2)\geq0$. For affine $F_t$ the second derivative is therefore nonpositive. If it vanishes, strict positivity of the smooth density of $\mu_{u_t}$ implies $A_t=0$ everywhere. Thus $\ddc\dot u_t=0$; the maximum principle gives that $\dot u_t$ is constant. The first differentiated equation then gives $G=-\dot b_t$, a constant. Conversely, if $G$ is constant, the solution metric is independent of $t$ and $b_t=b_0-tG$.
\end{proof}

\Needspace{6\baselineskip}
\begin{corollary}\label{cor:b-lipschitz}
For two smooth functions $F,G$, the normalization constants satisfy
\[
 -\sup_X(G-F)\leq b(G)-b(F)\leq-\inf_X(G-F).
\]
In particular, $|b(G)-b(F)|\leq\|G-F\|_\infty$. If $G-F$ is nonconstant, both displayed inequalities are strict.
\end{corollary}
\begin{proof}
Apply \eqref{eq:b-first} along $F_t=F+t(G-F)$ and integrate over $0\leq t\leq1$. Each $\mu_{u_t}$ is a probability measure with strictly positive smooth density. The integral of a nonconstant continuous function against it lies strictly between that function's minimum and maximum.
\end{proof}

Thus $b$ is a decreasing, concave, Lipschitz functional of the logarithmic density, with the translation rule $b(F+a)=b(F)-a$ for every real constant $a$. It extends uniquely as a Lipschitz functional from smooth to continuous $F$ by uniform approximation. This last statement concerns the scalar functional only; existence or regularity of limiting potentials requires further estimates.

\section{Normalization on the Hopf threefold}
\label{sec:hopf-normalization}

The Hopf metric of \cref{chap:hopf} makes the discrepancy from the closed-background normalization explicit: it is the energy of the solution along the elliptic fibers. In this section use that threefold, its projection $\pi:X\to\mathbb P^2$, the probability measure $\nu=\omega^3/6$ and the averaging operator $A$ of \cref{prop:hopf-fiber-average}. Write $\beta=\omega-\pi^*\omega_{\FS}$ for the rank-one form in that construction. The fiber derivative $u_w$ is independent of the base trivialization.

For $F\in C^\infty(X,\mathbb R)$, let $(u_F,b(F))$ be the solution of
\[
 (\omega+\ddc u_F)^3=\mathrm{e}^{F+b(F)}\omega^3,
 \qquad \omega+\ddc u_F>0,\qquad \int_Xu_F\,\mathrm{d}\nu=0.
\]
Define
\[
 D(F)=b(F)+\log\int_X\mathrm{e}^F\,\mathrm{d}\nu.
\]
The corresponding expression vanishes for a closed background by \eqref{eq:closed-b}.

\Needspace{13\baselineskip}
\begin{proposition}\label{prop:hopf-normalization-defect}\index{Hopf manifold!normalization defect}
For every smooth real $F$ on this threefold,
\begin{align}
 \mathrm{e}^{D(F)}&=1+\int_X|(u_F)_w|^2\,\mathrm{d}\nu,
 \label{eq:hopf-normalization-energy}\\
 \|u_F-Au_F\|_{L^2(X,\nu)}^2
 &\leq4\bigl(\mathrm{e}^{D(F)}-1\bigr).
 \label{eq:hopf-normalization-average}
\end{align}
Thus $D(F)\geq0$, with equality if and only if $F=\pi^*H$ for some $H\in C^\infty(\mathbb P^2,\mathbb R)$.
\end{proposition}
\begin{proof}
Integration of the equation and \eqref{eq:hopf-mass} give
\[
 \mathrm{e}^{D(F)}=\frac16\int_X(\omega+\ddc u_F)^3
                  =1+\int_X|(u_F)_w|^2\,\mathrm{d}\nu.
\]
Substituting the first equality into \eqref{eq:hopf-average-mass} proves \eqref{eq:hopf-normalization-average}.

If $D(F)=0$, smoothness gives $(u_F)_w=0$ everywhere. Since $u_F$ is real, it is constant on each connected fiber. It descends to a smooth real function $q$ on $\mathbb P^2$, with $\omega_{\FS}+\ddc q>0$. Since $\beta^2=0$ and $\omega=\pi^*\omega_{\FS}+\beta$, we have
\[
 (\omega+\ddc\pi^*q)^3
       =3\pi^*(\omega_{\FS}+\ddc q)^2\wedge\beta.
\]
Dividing by $\omega^3=3\pi^*\omega_{\FS}^2\wedge\beta$ and taking logarithms proves that $F$ is a smooth base pullback.

Conversely, suppose $F=\pi^*H$. Apply \cref{thm:tw-smooth} on $(\mathbb P^2,\omega_{\FS})$ to obtain a smooth real $q$ and a scalar $c$ with
\[
 \omega_{\FS}+\ddc q>0,\qquad
 (\omega_{\FS}+\ddc q)^2=\mathrm{e}^{H+c}\omega_{\FS}^2.
\]
Closedness of the base form and Stokes' theorem give $\mathrm{e}^c\int_{\mathbb P^2}\mathrm{e}^H\omega_{\FS}^2=1$. The preceding cube identity shows that $\pi^*q$, after subtraction of its $\nu$-mean, solves the equation on $X$ with scalar $c$. Uniqueness gives $b(F)=c$. Disintegration of $\nu$ now yields $D(F)=0$.
\end{proof}

In particular, if a smooth positive density $f$ satisfies $\int_Xf\,\mathrm{d}\nu=1$, the equation $(\omega+\ddc u)^3=\mathrm{e}^b f\omega^3$ has $b\geq0$, with equality precisely when $f$ comes from the base. Fixing the total integral of the density therefore does not determine the scalar on this non-closed background.

\begin{corollary}\label{cor:hopf-normalization-hessian}
Let $G\in C^\infty(X,\mathbb R)$, and let $v$ be the unique smooth real function satisfying
\[
 L_0v=G-\int_XG\,\mathrm{d}\nu,
 \qquad \int_Xv\,\mathrm{d}\nu=0.
\]
Then
\[
 \left.\frac{\mathrm{d}}{\mathrm{d}t}D(tG)\right|_{t=0}=0,
 \qquad
 \left.\frac{\mathrm{d}^2}{\mathrm{d}t^2}D(tG)\right|_{t=0}
       =2\int_X|v_w|^2\,\mathrm{d}\nu.
\]
The second derivative vanishes if and only if $G$ is a smooth base pullback.
\end{corollary}
\begin{proof}
For real $t$ near zero put $u_t=u_{tG}$ and $b_t=b(tG)$. By \cref{prop:smooth-dependence}, this is a smooth family, with $u_0=b_0=0$. Differentiating the logarithmic equation at zero gives $L_0\dot u_0=G+\dot b_0$. Since $\omega$ is Gauduchon, $\nu$ is the adjoint probability of $L_0$. Integration gives $\dot b_0=-\int_XG\,\mathrm{d}\nu$, so $\dot u_0=v$. This also proves existence and uniqueness of the displayed mean-zero solution.

Smooth dependence gives $(u_t)_w=t v_w+O(t^2)$ uniformly on $X$. Equation \eqref{eq:hopf-normalization-energy} therefore gives
\[
 \mathrm{e}^{D(tG)}=1+t^2\int_X|v_w|^2\,\mathrm{d}\nu+O(t^3).
\]
Taking logarithms proves the derivative formulas.

\Needspace{8\baselineskip}
If the second derivative vanishes, then $v=\pi^*q$ for a smooth real $q$ on the base. Direct wedge expansion gives
\[
 L_0(\pi^*q)=\pi^*\!\left(
       \frac{2\ddc q\wedge\omega_{\FS}}{\omega_{\FS}^2}\right).
\]
Hence $G=L_0v+\int_XG\,\mathrm{d}\nu$ is a base pullback. Conversely, when $G$ is a base pullback, the proposition gives $D(tG)=0$ for every real $t$, proving the converse.
\end{proof}

This is a second-variation formula at the background density. It does not determine global convexity of $D$: the second derivatives of $b$ and of the logarithmic integral have opposite signs.

\section{Prescribing a Gauduchon volume form}

We now return to a general connected compact complex manifold $X$ of dimension $n\geq2$. To preserve a Gauduchon Aeppli class, we vary the $(n-1)$st power of the metric by a $\partial$- and $\bar\partial$-exact term. The resulting scalar ansatz contains first derivatives of the potential as well as its complex Hessian.

\Needspace{22\baselineskip}
\begin{theorem}[Sz\'ekelyhidi--Tosatti--Weinkove]
\label{thm:stw-volume}\index{Gauduchon metric!prescribed volume}
Fix a Hermitian form $\gamma_0$, a Gauduchon form $\gamma$, and smooth real $F$. There is a unique pair $(u,b)\in C^\infty(X,\mathbb R)\times\mathbb R$ with $\sup_Xu=0$ such that the positive $(n-1,n-1)$-form
\[
 \Gamma_u=\gamma_0^{n-1}+\ddc u\wedge\gamma^{n-2}
 +\operatorname{Re}\left(
     \frac{\mathrm i}{2\pi}\partial u\wedge\bar\partial(\gamma^{n-2})
   \right)
\]
has a positive root $\eta_u^{n-1}=\Gamma_u$ satisfying
\[
 \eta_u^n=\mathrm{e}^{F+b}\gamma^n.
\]
\end{theorem}

Keep $\gamma_0$ and $\gamma$ fixed for the remainder of the chapter, and write $\mathcal D_\gamma u=\Gamma_u-\gamma_0^{n-1}$ for the linear part of the ansatz.

To see the preserved geometry, set
$\xi_u=\mathrm i\,\bar\partial u\wedge\gamma^{n-2}/(4\pi)$.
The Leibniz rule gives
$\Gamma_u-\gamma_0^{n-1}=\partial\xi_u+\bar\partial\overline{\xi_u}$.
If $\gamma_0$ is Gauduchon, so is $\eta_u$, and their $(n-1)$st powers have the same Aeppli class: the difference has the displayed exact form. Uniqueness here concerns the stated scalar ansatz; it does not assert uniqueness among all Gauduchon metrics in that Aeppli class.

\section{A uniform potential bound for the Gauduchon ansatz}
\label{sec:stw-czero}

Let $X$ be a connected compact complex manifold of dimension $n\geq2$. Fix a Hermitian form $\gamma_0$ and a Gauduchon form $\gamma$ on $X$. For smooth real $u$, let $\Gamma_u$ be the $(n-1,n-1)$-form in the scalar ansatz of \cref{thm:stw-volume}; when $\Gamma_u>0$, let $\eta_u$ be its positive root, $\eta_u^{n-1}=\Gamma_u$. We study $\eta_u^n=\mathrm{e}^{F+b}\gamma^n$ for prescribed smooth real $F$ and an unknown real scalar $b$.

We can prove the zeroth-order estimate in this equation without invoking its existence theorem. The argument uses a fixed linear Green operator and a local supporting-plane construction. The latter is the elementary mechanism behind the real Alexandrov--Bakelman--Pucci estimate. The supporting-plane construction gives a determinant estimate adapted to the Gauduchon ansatz.

\begin{proposition}\label{prop:stw-czero-reconstructed}
Fix $\gamma_0$ and $\gamma$ as in \cref{thm:stw-volume}. If a smooth pair $(u,b)$ satisfies that equation, with $\Gamma_u>0$ and $\sup_Xu=0$, then
\[
 |b|+\|u\|_{L^\infty(X)}\leq C.
\]
The constant depends only on the fixed geometry and an upper bound for $\|F\|_\infty$. It does not depend on derivatives of $F$ or on a bounded mass hypothesis.
\end{proposition}
\begin{proof}
We first bound the scalar. At a maximum of $u$, its first derivatives vanish and $\ddc u\leq0$, so $\Gamma_u\leq\gamma_0^{n-1}$. At a minimum the reverse inequality holds. The determinant of the positive cofactor matrix is monotone under the order of positive matrices; the positive-root volume is its $1/(n-1)$ power. Evaluating the equation at these two points gives
\[
 \inf_X\left(\log\frac{\gamma_0^n}{\gamma^n}-F\right)
 \leq b\leq
 \sup_X\left(\log\frac{\gamma_0^n}{\gamma^n}-F\right).
\]
This proves the scalar bound.

Write
\[
 \mathcal D_\gamma u
 =\ddc u\wedge\gamma^{n-2}
 +\operatorname{Re}\left(\frac{\mathrm i}{2\pi}
                \partial u\wedge\bar\partial\gamma^{n-2}\right)
\]
and consider the fixed real linear elliptic operator
\[
 Lu=\frac{\gamma\wedge\mathcal D_\gamma u}{\gamma^n}.
\]
Its principal part is $\ddc u\wedge\gamma^{n-1}/\gamma^n$, its other term is first order, and it has no zeroth-order term. Positivity of $\Gamma_u$ implies $Lu\geq-C_0$ for a fixed constant $C_0$.

We recall explicitly the standard linear Green fact used here. A fixed smooth real uniformly elliptic operator with no zeroth-order term on a connected compact manifold has a smooth strictly positive adjoint probability $\nu$. With the sign for which $L$ is nonpositive at a maximum, its normalized Green kernel can be written
\[
 \begin{gathered}
 v(x)-\int_Xv\,\mathrm{d}\nu
       =-\int_X K(x,y)Lv(y)\,\mathrm{d}\nu(y),\\
 K\geq0,\qquad \sup_x\int_XK(x,y)\,\mathrm{d}\nu(y)<+\infty.
 \end{gathered}
\]
Here the kernel is smooth off the diagonal and locally integrable. To obtain this version, construct the inverse on functions of zero $\nu$-mean by the Fredholm alternative. The local fundamental solution has a positive leading singularity after the displayed minus sign, so compactness bounds its negative part. Adding a constant to $K$ makes it nonnegative and does not change the formula, since $\int_XLv\,\mathrm{d}\nu=0$. The usual local Green estimates give the uniform integral bound. These are standard linear elliptic Green estimates; no nonlinear equation estimate is being used. The positive adjoint density follows from the maximum-principle argument already used in the proof of \cref{thm:gauduchon}.

Apply the formula to $u$ and use $Lu\geq-C_0$. At a maximum point, where $u=0$, it gives
\[
 -\int_Xu\,\mathrm{d}\nu\leq C_1.
\]
Because $u\leq0$ and $\nu$ has a fixed smooth positive density, this is a uniform $L^1$ bound in any fixed smooth volume form.

Let $M=\min_Xu$ and choose holomorphic coordinates on a ball centered at a minimum point. Compactness permits a fixed coordinate radius and uniform bounds for the metric coefficients and their derivatives; rescale the ball to the unit ball $B\subset\R^{2n}$. Write $x$ for its real coordinates, $D$ and $D^2$ for the Euclidean gradient and Hessian, $\mathrm{d}x$ for Euclidean volume, and $I$ for an identity matrix of the indicated size. Fix a small number $\varepsilon>0$, depending only on those bounds, and set
\[
 v(x)=u(x)+\varepsilon|x|^2.
\]
For any real vector $p$ with $|p|<\varepsilon/2$, the function $v(x)-p\cdot x$ has an interior minimum: its boundary values are at least $M+\varepsilon/2$, whereas its value at the center is $M$. At such a point,
\[
 Dv=p,\qquad D^2v\geq0,\qquad
 v\leq M+\varepsilon/2.
\]
Let $A\subset B$ be the set of all points satisfying
$|Dv|\leq\varepsilon/2$, $D^2v\geq0$, and $v\leq M+\varepsilon/2$.
Since $u\geq M$, this set lies in $\{|x|\leq1/\sqrt2\}$. It is closed there, hence compact in $B$, and its gradient image contains the ball of radius $\varepsilon/2$. The area formula for the smooth gradient map therefore gives
\begin{equation}\label{eq:stw-contact-area}
 |B_{\varepsilon/2}|
       \leq\int_A\det D^2v\,\mathrm{d}x.
\end{equation}
The Hessian is nonnegative on $A$, so no absolute value is needed in the determinant.

We bound this determinant pointwise on $A$. There,
\[
 |Du|\leq\frac52\varepsilon,
 \qquad D^2u\geq-2\varepsilon I.
\]
The complex Hessian has the corresponding uniform lower bound. Wedge with the fixed positive $\gamma^{n-2}$ and bound the first-order term using the gradient estimate. Let $Q_u$ be the matrix of $\Gamma_u$ in a fixed local frame of $(n-1,0)$-forms. Since $\gamma_0^{n-1}$ is strictly positive, choosing $\varepsilon$ small gives
\[
 Q_u\geq cI
\]
with a uniform $c>0$. The equation and the already proved bounds for $F+b$ bound $\det Q_u$ from above, because the determinant of this cofactor matrix is, up to a fixed positive frame factor, the $(n-1)$st power of $\eta_u^n$. All eigenvalues of $Q_u$ are therefore bounded above as well: their product is bounded and each of the other $n-1$ eigenvalues is at least $c$.

In particular $\gamma\wedge\Gamma_u/\gamma^n$ is uniformly bounded on $A$. Subtracting its fixed background term and its bounded first-order term bounds
\[
 \frac{\ddc u\wedge\gamma^{n-1}}{\gamma^n}
\]
from above. This is a positive fixed multiple of the $\gamma$-trace of the complex Hessian of $u$. The complex Hessian of $v$ is nonnegative on $A$, since the real Hessian is nonnegative. Metric comparison thus bounds its Euclidean complex trace from above. The real trace of $D^2v$ is four times that complex trace. It follows that
\[
 0\leq\det D^2v
 \leq\left(\frac{\operatorname{tr}_{\R}D^2v}{2n}\right)^{2n}
 \leq C_2\quad\hbox{on }A.
\]
Equation \eqref{eq:stw-contact-area} now gives $|A|\geq c_2>0$, where $c_2$ is independent of $u$. On $A$ we also have $u\leq v\leq M+\varepsilon/2$. Consequently the uniform $L^1$ bound gives
\[
 C_3\geq\int_A(-u)\,\mathrm{d}x
       \geq c_2(-M-\varepsilon/2).
\]
This bounds $-M$, and completes the proof.
\end{proof}

The estimate supplies the uniform potential bound needed in a continuity method for the Gauduchon ansatz. The second-order estimate and the rescaling argument below will supply the derivative bounds, retaining the first-order torsion terms.


\section{The adjoint equation for Gauduchon volume}
\label{sec:gauduchon-volume-adjoint}

Let $X$ be a connected compact complex manifold of dimension $n\geq2$. Fix a Hermitian form $\gamma_0$ and a Gauduchon form $\gamma$ on $X$. All potentials and test functions in this section are smooth real functions on $X$, and $u$ denotes a potential for the scalar ansatz of \cref{thm:stw-volume}. The linearization can be studied at any potential for which the ansatz is positive; this requires no existence theorem. It supplies the adjoint measure used in the continuity argument. The variation formulas for solutions follow once that argument is complete. They describe a second role for Gauduchon metrics, different from testing the ordinary Hermitian Monge--Amp\`ere equation.

Put $K=\gamma^{n-2}$ and define the real linear operator with values in $(n-1,n-1)$-forms by
\[
 \mathcal D_\gamma v
 =\ddc v\wedge K+
   \operatorname{Re}\left(\frac{\mathrm i}{2\pi}
              \partial v\wedge\bar\partial K\right).
\]
Thus $\Gamma_u=\gamma_0^{n-1}+\mathcal D_\gamma u$ and $\eta_u^{n-1}=\Gamma_u$. Fix a smooth positive volume form $\mathrm{d}V_0$. When differentiating solutions we use $\int_Xu\,\mathrm{d}V_0=0$ in place of the supremum normalization; this changes neither $\Gamma_u$ nor the scalar $b$.

\begin{proposition}\label{prop:gaud-volume-linearization}
On the open set $\Gamma_u>0$, the derivative of
\[
 \mathcal G(u)=\log\frac{\eta_u^n}{\gamma^n}
\]
is the real elliptic operator
\begin{equation}\label{eq:gaud-volume-linearization}
 \mathcal L_uv=
 \frac{n}{n-1}\frac{\eta_u\wedge\mathcal D_\gamma v}{\eta_u^n}.
\end{equation}
It has no zeroth-order term and its kernel consists of constants. There is a unique smooth positive function $\rho_u$, normalized by $\int_X\rho_u\eta_u^n=1$, such that the probability measure
\[
 \mathrm{d}\mu_u^\gamma=\rho_u\eta_u^n
\]
satisfies $\int_X\mathcal L_uv\,\mathrm{d}\mu_u^\gamma=0$ for every smooth real $v$.
\end{proposition}
\begin{proof}
For a smooth path of potentials, write $\eta=\eta_{u(t)}$ and $\Gamma=\Gamma_{u(t)}$, with a dot denoting differentiation in the real parameter $t$. Differentiating $\eta^{n-1}=\Gamma$ gives
$(n-1)\dot\eta\wedge\eta^{n-2}=\dot\Gamma$.
Wedge with $\eta$ and differentiate $\eta^n$. The result is
\[
 \frac{\mathrm{d}}{\mathrm{d}t}\log\eta^n
 =\frac{n}{n-1}\frac{\eta\wedge\dot\Gamma}{\eta^n},
\]
which proves \eqref{eq:gaud-volume-linearization}.

Only the first term of $\mathcal D_\gamma$ contributes to its principal symbol. For a nonzero complex covector $\zeta$, the form
\[
 \eta_u\wedge\mathrm i\zeta\wedge\overline\zeta
                       \wedge\gamma^{n-2}
\]
is strictly positive. Hence the principal symbol is positive definite, and compactness makes the smooth operator uniformly elliptic at each fixed $u$. The remaining term is first order, and constants are annihilated. The strong maximum principle gives $\ker\mathcal L_u=\mathbb R$.

The Fredholm and positivity arguments in the proof of \cref{thm:gauduchon} apply to this real elliptic operator with first-order terms and no zeroth-order term. Its index is zero by deformation to a Laplacian, so the adjoint for $\eta_u^n$ has a one-dimensional smooth kernel with a strictly positive generator. Normalizing its integral to one gives $\rho_u$ and uniqueness.
\end{proof}

The adjoint can be expressed as an equation for a positive $(1,1)$-form. For a smooth real $(1,1)$-form $A$, set
\[
 \mathcal A_\gamma(A)
 =\ddc A\wedge\gamma^{n-2}
 +\operatorname{Re}\left(\frac{\mathrm i}{2\pi}
              \partial A\wedge\bar\partial(\gamma^{n-2})\right).
\]
This is a real top-degree form.

\begin{proposition}\label{prop:gaud-volume-form-adjoint}
The multiplier in \cref{prop:gaud-volume-linearization} is characterized by
\begin{equation}\label{eq:gaud-volume-form-adjoint}
 \mathcal A_\gamma(\rho_u\eta_u)=0,
 \qquad \int_X\rho_u\eta_u^n=1.
\end{equation}
When $n=2$, its first equation is $\ddc(\rho_u\eta_u)=0$.
When $\eta_u=\gamma$, the multiplier is constant and
$\mathrm{d}\mu_u^\gamma=\gamma^n/\int_X\gamma^n$.
\end{proposition}
\begin{proof}
For a smooth real $(1,1)$-form $A$ and a smooth real function $v$ on $X$, we prove the integration formula
\begin{equation}\label{eq:gaud-volume-integration}
 \int_X A\wedge\mathcal D_\gamma v
 =\int_X v\,\mathcal A_\gamma(A).
\end{equation}
Write $c=\mathrm i/(2\pi)$. Since $A,K,v$ are real, expanding the real part gives
\[
 \mathcal D_\gamma v
 =c\,\partial\bar\partial v\wedge K
  +\frac c2\bigl(\partial v\wedge\bar\partial K
                   -\bar\partial v\wedge\partial K\bigr).
\]
Twice integrating the first term by parts gives
$\int_Xv\,c\partial\bar\partial(A\wedge K)$.
Integrating each first-order term once gives the correction
\[
 \int_Xv\,\frac c2
 \bigl(-\partial A\wedge\bar\partial K
       +\bar\partial A\wedge\partial K
       -2A\wedge\partial\bar\partial K\bigr).
\]
Finally,
\[
 \partial\bar\partial(A\wedge K)
 =\partial\bar\partial A\wedge K
   -\bar\partial A\wedge\partial K
   +\partial A\wedge\bar\partial K
   +A\wedge\partial\bar\partial K.
\]
Adding cancels the terms containing $A\wedge\partial\bar\partial K$ and leaves precisely $\mathcal A_\gamma(A)$. This proves \eqref{eq:gaud-volume-integration}. Apply it to $A=\rho_u\eta_u$ in \eqref{eq:gaud-volume-linearization}. It identifies the adjoint condition with \eqref{eq:gaud-volume-form-adjoint}.

For $n=2$, $K=1$, proving the first specialization. For the second, the product rule gives
\[
 \mathcal A_\gamma(\gamma)
 =\frac1{n-1}\ddc(\gamma^{n-1})=0.
\]
The unique normalized multiplier is consequently the constant $(\int_X\gamma^n)^{-1}$.
\end{proof}

For $n>2$ the equation for $\rho_u\eta_u$ contains the first derivatives of $\gamma^{n-2}$. The Gauduchon condition on $\eta_u$ alone is not the displayed adjoint equation. The fixed auxiliary metric $\gamma$ remains part of it.

\section{Normalization for the Gauduchon equation}

Let $X$ be a connected compact complex manifold of dimension $n\geq2$, with fixed Hermitian form $\gamma_0$ and Gauduchon form $\gamma$. Use the real linear form-valued operator $\mathcal D_\gamma$ of \cref{sec:gauduchon-volume-adjoint}, and put $\Gamma_u=\gamma_0^{n-1}+\mathcal D_\gamma u=\eta_u^{n-1}>0$ for smooth real potentials $u$. Thus $\eta_u$ is the positive root of the ansatz. Write $\mathcal G(u)=\log(\eta_u^n/\gamma^n)$, let $\mathcal L_u$ be its derivative from \cref{prop:gaud-volume-linearization}, and let $\mathrm{d}\mu_u^\gamma=\rho_u\eta_u^n$ be the corresponding normalized adjoint probability measure. Fix a smooth positive volume form $\mathrm{d}V_0$ and normalize potentials by $\int_Xu\,\mathrm{d}V_0=0$.

Let $b_\gamma(F)$ denote the scalar in \cref{thm:stw-volume}, with $\gamma_0$ and $\gamma$ fixed. It is defined for every smooth real $F$ by that existence theorem. We first justify differentiating it.

\begin{lemma}\label{lem:gaud-volume-smooth-dependence}
The mean-normalized pair $(u,b_\gamma(F))$ depends smoothly on smooth finite-dimensional families of $F$.
\end{lemma}
\begin{proof}
Fix an integer $k\geq0$ and $0<\alpha<1$. The positive-root map $\Gamma\mapsto\eta$ is smooth on positive $(n-1,n-1)$-forms: in a frame it is the inverse of the positive cofactor-matrix map. Thus $\mathcal G(u)-F-b$ is a smooth map into $C^{k,\alpha}(X,\mathbb R)$ on the open set of mean-zero $u\in C^{k+2,\alpha}(X,\mathbb R)$ with $\Gamma_u>0$, together with $b\in\mathbb R$ and $F\in C^{k,\alpha}(X,\mathbb R)$. Its derivative in $(u,b)$ is
\[
 (v,a)\longmapsto\mathcal L_uv-a,
 \qquad \int_Xv\,\mathrm{d}V_0=0.
\]
Given $h\in C^{k,\alpha}(X,\mathbb R)$, integrating against $\mu_u^\gamma$ forces $a=-\int_Xh\,\mathrm{d}\mu_u^\gamma$. The Fredholm alternative then solves for $v$, uniquely after the mean normalization. Schauder estimates give a bounded inverse. The implicit function theorem gives a local solution map around any given positive solution. This step uses only the adjoint construction. Once global existence and uniqueness in \cref{thm:stw-volume} are established, uniqueness joins these local maps. Arbitrarily high H\"older orders give the assertion for smooth data.
\end{proof}

\Needspace{7\baselineskip}
For a smooth family of prescribed functions $F_t\in C^\infty(X,\mathbb R)$, with $t$ in a real open interval, let $u_t$ be the mean-normalized solutions and let a dot denote differentiation in $t$. Choose any local frame for $(n-1,0)$-forms and let $Q_t$ be the positive Hermitian matrix of $\Gamma_{u_t}$ in the corresponding convention for positive $(n-1,n-1)$-forms. The expression
\[
 \operatorname{tr}\bigl((Q_t^{-1}\dot Q_t)^2\bigr)
\]
is independent of the chosen frame and is nonnegative: the endomorphism is similar to the Hermitian matrix $Q_t^{-1/2}\dot Q_t Q_t^{-1/2}$.

\Needspace{18\baselineskip}
\begin{theorem}\label{thm:gaud-volume-normalization-derivatives}
For a smooth family $F_t$ and its solutions,
\begin{align*}
 \dot b_\gamma(F_t)&=-\int_X\dot F_t\,\mathrm{d}\mu_{u_t}^\gamma,\\
 \ddot b_\gamma(F_t)&=-\int_X\ddot F_t\,\mathrm{d}\mu_{u_t}^\gamma
 -\frac1{n-1}\int_X
        \operatorname{tr}\bigl((Q_t^{-1}\dot Q_t)^2\bigr)
                          \,\mathrm{d}\mu_{u_t}^\gamma.
\end{align*}
In particular $b_\gamma$ is concave, and along $F_t=F_0+tG$ its second derivative is strictly negative unless $G$ is constant. For smooth real functions $F,G$ on $X$ and a real constant $a$, moreover,
\[
 b_\gamma(F+a)=b_\gamma(F)-a,
 \qquad |b_\gamma(G)-b_\gamma(F)|\leq\|G-F\|_\infty.
\]
\end{theorem}
\begin{proof}
The determinant of the cofactor matrix of a positive $n\times n$ matrix $H$ is $(\det H)^{n-1}$. Consequently, up to a term independent of $u$ coming from the fixed frame,
\[
 \mathcal G(u)=\frac1{n-1}\log\det Q.
\]
Differentiate $\mathcal G(u_t)=F_t+b_\gamma(F_t)$ once. Integration against the adjoint probability gives the first formula.

Because $\Gamma_u$ is affine in $u$, its second derivative along the family is $\mathcal D_\gamma\ddot u_t$. The matrix derivative identity for $\log\det$ therefore gives
\[
 \mathcal L_{u_t}\ddot u_t
 -\frac1{n-1}\operatorname{tr}\bigl((Q_t^{-1}\dot Q_t)^2\bigr)
 =\ddot F_t+\ddot b_\gamma(F_t).
\]
Integrate at this fixed parameter against $\mu_{u_t}^\gamma$ to obtain the second formula. No derivative of the measure is needed.

For affine $F_t$, positivity of the density of $\mu_{u_t}^\gamma$ proves concavity. Equality in the second derivative forces $\dot\Gamma_{u_t}=0$, hence $\mathcal L_{u_t}\dot u_t=0$. The kernel statement makes $\dot u_t$ constant, so the first differentiated equation forces $G$ to be constant. Conversely a constant change of $F$ changes only $b$ by the opposite constant. Finally integrate the first derivative along the line from $F$ to $G$; it lies between $-\sup_X(G-F)$ and $-\inf_X(G-F)$, proving the stated Lipschitz bound.
\end{proof}

This calculation is the normalization-concavity method already used for the ordinary Hermitian equation, with the determinant exponent $1/(n-1)$ and the adjoint operator appropriate to the Gauduchon ansatz. It is a consequence of the smooth equation, not an independent proof of its existence.

\Needspace{9\baselineskip}
\begin{example}\label{ex:gaud-volume-surface}
If $n=2$ and $\gamma_0$ is Gauduchon, then $\eta_u=\gamma_0+\ddc u$ is also Gauduchon. Stokes gives
\[
 \int_X\eta_u^2=\int_X\gamma_0^2,
\]
since the linear term pairs $u$ with $\ddc\gamma_0=0$ and the quadratic term integrates to zero. Therefore
\[
 b_\gamma(F)=\log\int_X\gamma_0^2-
                   \log\int_X\mathrm{e}^F\gamma^2,
 \qquad
 \mathrm{d}\mu_u^\gamma=\frac{\mathrm{e}^F\gamma^2}{\int_X\mathrm{e}^F\gamma^2}.
\]
The adjoint equation gives the same measure because $\rho_u$ is constant. Along $F+tG$, the second derivative is minus the variance of $G$ under this probability. This verifies the normalization formula directly in a case with conserved mass.
\end{example}

\Needspace{7\baselineskip}
The Lipschitz bound extends the scalar $b_\gamma$ uniquely to continuous $F$. It does not construct a weak positive $(n-1,n-1)$-form solving the equation. Such a theory would need a definition of the nonlinear root-volume measure for singular forms in the Aeppli ansatz, compactness of the corresponding potentials, and a stability theorem for the adjoint equation. These are separate questions; ordinary weak solutions for a smooth real $(1,1)$-background $\theta$ from the preceding sections do not supply them.


\Needspace{11\baselineskip}
\section{The matrix structure of the Gauduchon equation}
\label{sec:stw-matrix-structure}

The Hodge star converts the positive $(n-1,n-1)$-form $\Gamma_u$ into a Hermitian matrix whose logarithmic determinant is a concave function of the complex Hessian. Two features will control the largest eigenvalue in the second-order estimate: the positive background gives coercivity, and the torsion coefficients omit exactly those gradient indices that would otherwise carry the largest inverse weight. We derive both properties before differentiating the equation.

For the matrix and derivative calculations, set $v=u/(2\pi)$, so that $\ddc u=\mathrm i\partial\bar\partial v$. All covariant derivatives are taken with respect to the Chern connection $\nabla$ of $\gamma$. In holomorphic coordinates $z^1,\ldots,z^n$, write $v_i=\partial_i v$ and $v_{\bar j}=\partial_{\bar j}v$; indices range from $1$ to $n$, and $\delta_{ij}$ is the Kronecker symbol. Write $\gamma=\mathrm i\gamma_{a\bar b}\,\mathrm{d}z^a\wedge\mathrm{d}\bar z^b$, with repeated indices summed. We denote by $\star_\gamma$ the Hodge star of $\gamma$, and by $\det_\gamma$ the determinant of the Hermitian matrix of a real $(1,1)$-form relative to $\gamma$. Write the torsion of $\nabla$ as $T_{ab}^{c}$ and lower the last index with $\gamma$. Thus
\[
 T_{ab\bar c}=\partial_a\gamma_{b\bar c}-\partial_b\gamma_{a\bar c},
 \qquad T_{ab\bar c}=-T_{ba\bar c}.
\]
Let $Q_0=\star_\gamma\gamma_0^{n-1}/(n-1)!$ and
\[
 Z=\frac1{(n-1)!}\star_\gamma
       \operatorname{Re}\bigl(\mathrm i\partial v\wedge
                    \bar\partial\gamma^{n-2}\bigr).
\]
Here the real $(1,1)$-forms on the right are identified with their Hermitian coefficient matrices relative to $\gamma$. Set
\[
 \chi=(\operatorname{tr}_\gamma Q_0)\gamma-(n-1)Q_0,
 \quad W=(\operatorname{tr}_\gamma Z)\gamma-(n-1)Z,
 \quad g=\chi+\mathrm i\partial\bar\partial v+W.
\]
The linear map $P_\gamma(A)=((\operatorname{tr}_\gamma A)\gamma-A)/(n-1)$ on real $(1,1)$-forms, or Hermitian matrices, is invertible, and the equation of \cref{thm:stw-volume} becomes
\begin{equation}\label{eq:stw-cofactor-matrix-scaling}
 \begin{aligned}
 Q_u&:=P_\gamma(g)=Q_0+P_\gamma(\mathrm i\partial\bar\partial v)+Z>0,\\
 \log\det_\gamma Q_u&=h=(n-1)(F+b).
 \end{aligned}
\end{equation}
Here $Q_u=\star_\gamma\Gamma_u/(n-1)!$ is the coefficient matrix of $\Gamma_u$. The tensor $g$ need not be positive. To check the Hessian term, diagonalize $\mathrm i\partial\bar\partial v$ in a unitary frame: the coefficient of $\star_\gamma(\mathrm i\partial\bar\partial v\wedge\gamma^{n-2})/(n-1)!$ in each diagonal position is the sum of the other diagonal entries divided by $n-1$. If $\eta_u$ has eigenvalues $a_i$ relative to $\gamma$, then $Q_u$ has eigenvalues $\prod_{j\ne i}a_j$, whose product is $(\prod_i a_i)^{n-1}$. This proves the determinant identity.

\Needspace{20\baselineskip}
\begin{lemma}\label{lem:stw-torsion-coefficients}
At a point in a unitary frame,
\begin{align}
 2(n-1)Z_{i\bar j}
 ={}&\delta_{ij}\sum_{p,k}v_p\overline{T_{pk\bar k}}
       -v_i\sum_k\overline{T_{jk\bar k}}
       -\sum_kv_k\overline{T_{kj\bar i}}\notag\\
 &+\delta_{ij}\sum_{p,k}v_{\bar p}T_{pk\bar k}
       -v_{\bar j}\sum_kT_{ik\bar k}
       -\sum_kv_{\bar k}T_{ki\bar j}.
 \label{eq:stw-torsion-six-terms}
\end{align}
\Needspace{9\baselineskip}
In particular $Z_{i\bar j}$ contains no $v_j$ or $v_{\bar i}$. Moreover
\begin{equation}\label{eq:stw-diagonal-torsion}
 2(n-1)Z_{i\bar i}
   =\sum_{p\ne i,\ k\ne i}
        \bigl(v_p\overline{T_{pk\bar k}}+v_{\bar p}T_{pk\bar k}\bigr).
\end{equation}
Covariant differentiation preserves these coefficient identities. For example $\nabla_1 Z_{1\bar1}$ contains no pure Hessian $v_{11}$ or its conjugate, and $\nabla_{\bar1}\nabla_1 Z_{1\bar1}$ contains no $v_{11\bar1}$ or its conjugate, nor a pure $v_{11}$ or its conjugate.
\end{lemma}
\begin{proof}
For $n=2$, the defining expression of $Z$ is zero, and so are the formulas below, because the generalized Kronecker symbol with three indices vanishes in two dimensions. Suppose $n\geq3$ and use coordinates that are unitary at the point. Direct antisymmetrization gives
\[
 \bar\partial\gamma
   =-\frac{\mathrm i}{2}\sum_{a,b,c}
        \overline{T_{bc\bar a}}\,\mathrm{d}z^a\wedge \mathrm{d}\bar z^b\wedge \mathrm{d}\bar z^c.
\]
We use $\bar\partial\gamma^{n-2}=(n-2)\bar\partial\gamma\wedge\gamma^{n-3}$. To write the contraction explicitly, put
\[
 \delta^{paj}_{bci}
   =\det\begin{pmatrix}
     \delta_{pb}&\delta_{pc}&\delta_{pi}\\
     \delta_{ab}&\delta_{ac}&\delta_{ai}\\
     \delta_{jb}&\delta_{jc}&\delta_{ji}
   \end{pmatrix}.
\]
Let $A=\star_\gamma(\mathrm i\partial v\wedge\bar\partial\gamma^{n-2})/(n-1)!$. Expanding the $n-3$ remaining copies of $\gamma$ leaves the three complementary indices and gives
\[
 A_{i\bar j}=\frac1{2(n-1)}
    \sum_{p,a,b,c}v_p\overline{T_{bc\bar a}}\,
                         \delta^{paj}_{bci}.
\]
Indeed, the factor is $(n-2)(n-3)!/[2(n-1)!]$, and the determinant records the signs of the remaining holomorphic and antiholomorphic indices. Expanding its six terms and using $T_{bc\bar a}=-T_{cb\bar a}$, we obtain
\[
 A_{i\bar j}=\frac1{n-1}\left(
   \delta_{ij}\sum_{p,k}v_p\overline{T_{pk\bar k}}
   -v_i\sum_k\overline{T_{jk\bar k}}
   -\sum_kv_k\overline{T_{kj\bar i}}\right).
\]
Since $Z=(A+A^*)/2$, where $A^*$ denotes the Hermitian adjoint matrix, this gives \eqref{eq:stw-torsion-six-terms}.

For $i\ne j$, the coefficient of $v_j$ in this expression vanishes because $T_{jj\bar i}=0$, and the coefficient of $v_{\bar i}$ vanishes in the conjugate expression. For $i=j$, cancelling the terms with $p=i$ and $k=i$ gives \eqref{eq:stw-diagonal-torsion}. These are tensor identities valid in every unitary frame. Since the Chern connection preserves the metric, they can be differentiated in a unitary frame that is parallel at the point. For instance, the first derivative of a summand in \eqref{eq:stw-diagonal-torsion} is
\[
 (\nabla_1v_p)\overline{T_{pk\bar k}}
       +v_p\nabla_1\overline{T_{pk\bar k}},\qquad p,k\ne1,
\]
together with its conjugate counterpart. A further derivative in the direction $\bar1$ still has $p\ne1$. This gives the stated exclusions, including the terms in which a derivative falls on a coefficient.
\end{proof}

The second-order estimate uses these exclusions; an arbitrary first-order perturbation could not be substituted for $Z$.

\Needspace{22\baselineskip}
\begin{lemma}\label{lem:stw-logdet-coercivity}
Diagonalize $g$ at a point and order its eigenvalues as $\lambda_1\geq\cdots\geq\lambda_n$. Put
\[
 \mu_i=\frac{\sum_{j\ne i}\lambda_j}{n-1}>0,
 \quad t_i=\mu_i^{-1},\quad
 f_i=\frac1{n-1}\sum_{j\ne i}t_j,
 \quad \mathcal F=\sum_i f_i=\sum_i t_i.
\]
Then $\mu_1\leq\cdots\leq\mu_n$, $f_1\leq\cdots\leq f_n$, and
\begin{gather}
 \sum_i f_i\lambda_i=n,\qquad
 \frac{\mathcal F}{n(n-1)}\leq f_k\leq\frac{\mathcal F}{n-1}
       \quad(k>1),\qquad t_i\leq(n-1)f_1\quad(i>1),
 \label{eq:stw-weight-comparability}\\
 f_n\geq\mathcal F/n.
 \notag
\end{gather}
Fix $\tau,D>0$ and assume $Q_0\geq\tau\gamma$ and $\det P_\gamma(g)\leq D$. For $\lambda_1$ larger than a constant depending only on $n,\tau,D$,
\begin{equation}\label{eq:stw-direct-coercivity}
 \sum_i f_i(\chi_{i\bar i}-\lambda_i)
                \geq\frac\tau2\mathcal F.
\end{equation}
There is also a uniform positive lower bound for $\mathcal F$ depending only on $n,D$.
\end{lemma}
\begin{proof}
The ordering follows from the formula for $\mu_i$, and then from that for $f_i$. The identity $\sum_if_i\lambda_i=n$, which reflects the relation $\log\det P_\gamma(sg)=\log\det P_\gamma(g)+n\log s$, can be checked directly:
\[
 \sum_i f_i\lambda_i
   =\sum_j\frac{t_j}{n-1}\sum_{i\ne j}\lambda_i
   =\sum_jt_j\mu_j=n.
\]
For $k>1$, the sum defining $f_k$ contains the largest weight $t_1\geq\mathcal F/n$. The upper bound for $f_k$ and the estimate for $t_i$ are immediate, and the largest weight $f_n$ is at least the average $\mathcal F/n$.

The map $P_\gamma$ is self-adjoint for the matrix trace pairing. Consequently
\[
 \sum_i f_i\chi_{i\bar i}
   =\sum_i t_i(Q_0)_{i\bar i}\geq\tau\mathcal F.
\]
Moreover $\sum_i\lambda_i=\sum_i\mu_i>\lambda_1>0$, where the strict inequality follows from $\mu_1>0$. Hence $\mu_n>\lambda_1/n$. Since the other $n-1$ numbers $\mu_i$ are at least $\mu_1$,
\[
 \mu_1^{n-1}\mu_n\leq\prod_i\mu_i\leq D,
 \qquad
 \mathcal F\geq\mu_1^{-1}
             \geq\left(\frac{\lambda_1}{nD}\right)^{1/(n-1)}.
\]
For $\lambda_1$ sufficiently large this is at least $2n/\tau$, so $n\leq\tau\mathcal F/2$. Subtracting $\sum_i f_i\lambda_i=n$ therefore proves \eqref{eq:stw-direct-coercivity}. Finally, the arithmetic--geometric mean inequality applied to the $t_i$ gives $\mathcal F\geq nD^{-1/n}$.
\end{proof}

\Needspace{8\baselineskip}
The concavity of the logarithmic determinant will supply a positive term with an explicit formula. For the matrix function $\Phi(g)=\log\det P_\gamma(g)$, the negative second differential in a complex matrix direction $B$ and its adjoint is
\begin{equation}\label{eq:stw-concavity-square}
 -D^2\Phi_g(B,B^*)
   =\operatorname{tr}\bigl(P_\gamma(g)^{-1}P_\gamma(B)
                   P_\gamma(g)^{-1}P_\gamma(B)^*\bigr)\geq0.
\end{equation}
This follows by differentiating the inverse matrix. In a unitary frame diagonalizing $g$, the off-diagonal terms of this expression include
\[
 \frac{|B_{k\bar1}|^2}{(n-1)^2\mu_k\mu_1}.
\]
When $\lambda_k<\lambda_1$, that coefficient is exactly
\[
 \frac{f_k-f_1}{\lambda_1-\lambda_k}
       =\frac1{(n-1)^2\mu_k\mu_1}.
\]
The expression on the right also makes sense when the two eigenvalues coincide. This positive square and the weights $f_i$ enter the largest-eigenvalue estimate in the next section.


\section{The largest-eigenvalue estimate}
\label{sec:stw-second-order-proof}

The estimate bounds the complex Hessian in terms of the gradient squared. At a maximum of a function of the largest eigenvalue, the gradient and the potential, the missing torsion indices allow the concavity square to absorb third-derivative terms. Coercivity closes the estimate when the smallest eigenvalue is controlled; in the other case its negative size supplies a positive Hessian square.

Plain subscripts on scalar functions and comma subscripts on matrix coefficients denote coordinate derivatives. The coefficients of the Chern connection are written $\Gamma^q_{ik}$. Choose $\tau>0$ with $Q_0\geq\tau\gamma$.

\Needspace{16\baselineskip}
\begin{theorem}\label{thm:stw-second-order-reconstructed}
Fix smooth Hermitian forms $\gamma_0,\gamma$ on a compact complex manifold of dimension $n\geq2$. For a smooth real function $h$ on $X$, suppose that $v$ is a smooth solution of
\[
 Q=Q_0+P_\gamma(\mathrm i\partial\bar\partial v)+Z(\partial v)>0,
 \qquad \log\det_\gamma Q=h,
\]
with the tensors of \cref{sec:stw-matrix-structure}. If $\|v\|_{C^0}$ and $\|h\|_{C^2}$ are bounded, then
\[
 \sup_X|\mathrm i\partial\bar\partial v|_\gamma
       \leq C\bigl(1+\sup_X|\partial v|_\gamma^2\bigr),
\]
where $C$ depends only on these bounds and the fixed smooth data.
\end{theorem}
\begin{proof}
Adding a constant, arrange $\sup v=0$ and write $-M\leq v\leq0$. Set
\[
 K=1+\sup_X|\partial v|_\gamma^2,\qquad
 \phi(s)=-\tfrac12\log(1-s/(2K)),\qquad
 \psi(t)=D_1\mathrm{e}^{-D_2t}.
\]
The positive constants $D_1,D_2$ will be chosen below using only the fixed data. On $0\leq s\leq K$,
\[
 \frac1{4K}\leq\phi'\leq\frac1{2K},\qquad
 \phi''=2(\phi')^2,\qquad 0\leq\phi\leq\tfrac12\log2.
\]
Let $\lambda_1$ be the largest eigenvalue of $g=\chi+\mathrm i\partial\bar\partial v+W$ relative to $\gamma$. It is positive because $\operatorname{tr}_\gamma g=\operatorname{tr}_\gamma Q>0$. Take a maximum of
\[
 H=\log\lambda_1+\phi(|\partial v|_\gamma^2)+\psi(v).
\]
It suffices to bound $\lambda_1/K$ at this point: the other two terms are uniformly bounded once $D_1,D_2$ are fixed. We may assume that this ratio exceeds any fixed constant specified in the course of the argument.

Choose holomorphic coordinates unitary at the maximum, with $g$ diagonal there and eigenvalues ordered $\lambda_1\geq\cdots\geq\lambda_n$. All estimates below are at this point. Repeated indices are summed only when a sum is displayed. Put $\lambda=\lambda_1$ and use the weights $t_i,f_k$ and their sum $\mathcal F$ from \cref{lem:stw-logdet-coercivity}. Since $\sum_i\lambda_i>0$ and the eigenvalues are ordered,
\[
 |\lambda_i|\leq(n-1)\lambda,\qquad
 |v_{i\bar j}|\leq C\lambda
       \quad\hbox{if }\lambda\geq K.
\]
The constant $C$ may change and is independent of $K$. Until $D_1,D_2$ are fixed below, it is also independent of them. Pure derivatives $v_{ij}$ are not included in this bound.

\paragraph{Differentiating the largest eigenvalue.}
We make its use at a multiple eigenvalue precise. Choose a smooth local unitary frame parallel for the Chern connection at the point, with first vector an eigenvector for $\lambda$. In this frame let $B$ have eigenvalues $0,\epsilon_0,\ldots,\epsilon_0$, where $\epsilon_0>0$ is fixed and small. Then $B\geq0$, $\nabla B=0$ at the point, and its second derivatives are bounded by fixed background constants. The largest eigenvalue $\widetilde\lambda$ of $g-B$ is simple near the point, is at most $\lambda_1(g)$, and equals $\lambda$ there. Hence the corresponding smooth function $\widetilde H$ has a local maximum there. Only the largest eigenvalue needs to be simple.

For a self-adjoint matrix family with a simple top eigenvalue, differentiate its normalized eigenvector equation. Its component in the $p$th eigenspace is the corresponding derivative of the matrix divided by the eigenvalue gap. Differentiating the eigenvalue once more therefore gives
\begin{align*}
 a_k:=\widetilde\lambda_k&=\nabla_k g_{1\bar1},\\
 \widetilde\lambda_{k\bar k}
 &=\nabla_{\bar k}\nabla_k g_{1\bar1}+O(1)
   +\sum_{p>1}\frac{|\nabla_k g_{1\bar p}|^2+
                            |\nabla_k g_{p\bar1}|^2}
                        {\lambda-\lambda_p+\epsilon_0}.
\end{align*}
\Needspace{11\baselineskip}
This formula remains valid if some lower eigenvalues coincide: sum over any orthonormal basis of their eigenspaces. The gaps are positive and at most $n\lambda+\epsilon_0$. In holomorphic coordinates, $a_k=g_{1\bar1,k}+O(\lambda)$, while converting the second covariant derivative to partial derivatives introduces at most
$C\sum_p|g_{1\bar p,k}|+C\lambda$. Convert the off-diagonal covariant derivatives in the positive squares as well; their errors are $O(\lambda)$. Young's inequality absorbs the linear off-diagonal terms. Thus, after decreasing a fixed positive coefficient,
\begin{equation}\label{eq:stw-eigenvalue-basic}
 \widetilde\lambda_{k\bar k}
 \geq g_{1\bar1,k\bar k}
    +\frac1{8n\lambda}\sum_{p>1}
          (|g_{1\bar p,k}|^2+|g_{p\bar1,k}|^2)
    -C(|g_{1\bar1,k}|+\lambda).
\end{equation}
The precise positive constant will not be used.

Write $V_k=\partial_k(|\partial v|_\gamma^2)$. The first derivative of $\widetilde H$ gives
\begin{equation}\label{eq:stw-maximum-first}
 \frac{a_k}{\lambda}=-\phi' V_k-\psi'v_k,
 \qquad
 |V_k|\leq C\sqrt K\sum_p(|v_{pk}|+|v_{p\bar k}|)+CK.
\end{equation}
For the second derivative, the derivatives of the fixed metric in $|\partial v|^2$ cost only a small fraction of its positive Hessian squares and a constant. Therefore
\begin{align}
 \widetilde H_{k\bar k}\geq{}&
  \frac{\widetilde\lambda_{k\bar k}}\lambda
       -\frac{|a_k|^2}{\lambda^2}
       +\phi'\sum_p(v_pv_{\bar p k\bar k}+v_{\bar p}v_{pk\bar k})
       +\frac1{5K}\sum_p(|v_{pk}|^2+|v_{p\bar k}|^2)\notag\\
 &+\phi''|V_k|^2+\psi''|v_k|^2+\psi'v_{k\bar k}-C.
 \label{eq:stw-maximum-second-preliminary}
\end{align}
The third term follows directly by differentiating $\sum\gamma^{p\bar q}v_pv_{\bar q}$ twice and applying $ab\leq\varepsilon a^2+C_\varepsilon b^2$ to each metric-derivative term.

\Needspace{10\baselineskip}
\paragraph{The third derivatives and the torsion term.}
We provide the estimates that remove third derivatives from this inequality. Define
\begin{gather*}
 E=\sum_{k,p} f_k(|v_{pk}|^2+|v_{p\bar k}|^2),\qquad
 S=\sum_{k,p} f_k(|v_{pk}|+|v_{p\bar k}|),\\
 \mathcal C=-D^2\Phi_g(\nabla_1g,(\nabla_1g)^*)\geq0.
\end{gather*}
\Needspace{4\baselineskip}
The first differentiated equation is
\begin{equation}\label{eq:stw-first-differentiated}
 \sum_k f_k\nabla_p g_{k\bar k}=h_p.
\end{equation}
Since $P_\gamma(W)=Z$ and the connection preserves $\gamma$, one has the exact contractions
\begin{equation}\label{eq:stw-contraction-derivatives}
 \sum_k f_k\nabla_p W_{k\bar k}=\sum_i t_i\nabla_p Z_{i\bar i},
 \qquad
 \sum_k f_k\nabla_{\bar1}\nabla_1W_{k\bar k}
       =\sum_i t_i\nabla_{\bar1}\nabla_1Z_{i\bar i}.
\end{equation}
The weights in these identities are evaluated at the point; they are not differentiated.

Here is how the large weights are controlled. The only weight $t_i$ not bounded by $(n-1)f_1$ is $t_1$. In $Z_{1\bar1}$, every gradient index differs from $1$. Thus a pure Hessian in $\nabla_pZ_{1\bar1}$ has an index $q>1$, and can be weighted by $f_q$, which is comparable to $\mathcal F$. Pure derivatives commute in these coordinates. A mixed Hessian can similarly be placed in column $q$ using conjugate symmetry. All other $t_i$ are bounded by a constant times every $f_k$. Derivatives of coefficients cost $C\sqrt K\mathcal F$. Applying these observations to the first identity, and then to \eqref{eq:stw-first-differentiated}, gives, for each $p$,
\begin{equation}\label{eq:stw-weighted-third-bound}
 \left|\sum_k f_kv_{k\bar k p}\right|
       +\left|\sum_k f_kv_{k\bar k\bar p}\right|
       \leq C(S+\sqrt K\mathcal F).
\end{equation}
The connection terms in the differentiated equation contribute weighted mixed Hessians and $C\sqrt K\mathcal F$, already on the right. Bounded derivatives of $h$ are absorbed using the positive lower bound for $\mathcal F$.

For clarity we also track the second derivative of the equation through the eigenvalue estimate. Partial derivatives of $v$ commute, so
\[
 g_{1\bar1,k\bar k}
   =g_{k\bar k,1\bar1}+W_{1\bar1,k\bar k}
          -W_{k\bar k,1\bar1}+O(1).
\]
Twice differentiating the equation covariantly gives
\[
 \sum_k f_k\nabla_{\bar1}\nabla_1g_{k\bar k}
       =h_{1\bar1}+\mathcal C.
\]
The conversion to partial derivatives uses
\[
 \nabla_{\bar1}\nabla_1g_{k\bar k}
       =g_{k\bar k,1\bar1}
           -2\operatorname{Re}\sum_q\overline{\Gamma^q_{1k}}g_{k\bar q,1}
           +O(\lambda).
\]
Rewrite the derivative of $g$ in this term using
$g_{k\bar q,1}=g_{1\bar q,k}+W_{k\bar q,1}-W_{1\bar q,k}+O(1)$ and its conjugate. The first part is absorbed by the off-diagonal squares in \eqref{eq:stw-eigenvalue-basic}, with cost $C\lambda\mathcal F+C\sum_k f_k|g_{1\bar1,k}|$. The $W_{1\bar q,k}$ part is bounded by
$C\lambda\mathcal F+C\sum_{k,p}f_k|v_{pk}|$. For the remaining part use
\[
 2\operatorname{Re}\sum_q\overline{\Gamma^q_{1k}}W_{k\bar q,1}
       =W_{k\bar k,1\bar1}
          -\nabla_{\bar1}\nabla_1W_{k\bar k}+O(\sqrt K).
\]
It cancels the partial second derivative $-W_{k\bar k,1\bar1}$ above. We obtain
\begin{align}
 \sum_k f_k\widetilde\lambda_{k\bar k}
 \geq{}&\mathcal C+\sum_k f_k
       (W_{1\bar1,k\bar k}-\nabla_{\bar1}\nabla_1W_{k\bar k})\notag\\
 &-C\left(\sum_k f_k|g_{1\bar1,k}|
                +\sum_{k,p}f_k|v_{pk}|+\lambda\mathcal F\right).
 \label{eq:stw-eigenvalue-torsion-remainder}
\end{align}

The second contraction in \eqref{eq:stw-contraction-derivatives} bounds its negative torsion term. In the $t_1$ part the possible third derivatives are $v_{q1\bar1}$ and their conjugates with $q>1$; all other parts have weights bounded by $Cf_1$. Commute partial derivatives and use
\[
 v_{1\bar1,k}=g_{1\bar1,k}-\chi_{1\bar1,k}-W_{1\bar1,k}.
\]
The exclusions in \cref{lem:stw-torsion-coefficients} ensure that any remaining pure $v_{11}$ carries only the small weight $Cf_1$. Derivatives of the connection contribute second derivatives with the same weighted bounds. It follows that
\[
 \left|\sum_k f_k\nabla_{\bar1}\nabla_1W_{k\bar k}\right|
 \leq C\left(\sum_k f_k|g_{1\bar1,k}|
             +\sum_{k,p}f_k|v_{pk}|+\lambda\mathcal F\right).
\]
For the positive torsion term write $W_{1\bar1}=\sum_p(c^pv_p+\overline{c^p}v_{\bar p})$ with fixed smooth coefficients. Two derivatives give $c^p\sum_k f_kv_{pk\bar k}$ and its conjugate, plus terms bounded by $C(S+\sqrt K\mathcal F)$. Equation \eqref{eq:stw-weighted-third-bound} bounds the former. In this estimate the mixed Hessians can be bounded by $C\lambda$, giving
\[
 \left|\sum_k f_k W_{1\bar1,k\bar k}\right|
       \leq C\left(\sum_{k,p}f_k|v_{pk}|+\lambda\mathcal F\right).
\]
\Needspace{12\baselineskip}
Substitute these estimates in \eqref{eq:stw-eigenvalue-torsion-remainder}, then in \eqref{eq:stw-maximum-second-preliminary}. The third derivatives in the latter are bounded by \eqref{eq:stw-weighted-third-bound} and $\phi'|\partial v|\leq C/\sqrt K$. The resulting linear Hessian terms are absorbed by a small part of $E/K$, with cost $C\mathcal F$. Thus
\begin{equation}\label{eq:stw-key-maximum-inequality}
 \begin{split}
 0\geq{}&\frac{\mathcal C}{\lambda}
       -\sum_k f_k\frac{|a_k|^2}{\lambda^2}
       +\frac{E}{6K}+\phi''\sum_k f_k|V_k|^2
       +\psi''\sum_k f_k|v_k|^2+\psi'\sum_k f_kv_{k\bar k}\\
 &-C\left(\sum_k f_k\frac{|g_{1\bar1,k}|}{\lambda}+\mathcal F\right).
 \end{split}
\end{equation}

\paragraph{Cancellation when the smallest eigenvalue is controlled.}
Fix $0<\delta<1/2$, to be chosen after $D_1,D_2$. First suppose $\lambda_n\geq-\delta\lambda$. Put $I=\{k:f_k>\delta^{-1}f_1\}$, so $1\notin I$. Equation \eqref{eq:stw-maximum-first} and $\phi''=2(\phi')^2$ give
\begin{align*}
 -\sum_{k\notin I}f_k|a_k|^2/\lambda^2
 &\geq-\phi''\sum_{k\notin I}f_k|V_k|^2
              -2(\psi')^2\delta^{-1}f_1K,\\
 -2\delta\sum_{k\in I}f_k|a_k|^2/\lambda^2
 &\geq-2\delta\phi''\sum_{k\in I}f_k|V_k|^2
              -4\delta(\psi')^2\sum_{k\in I}f_k|v_k|^2.
\end{align*}
\Needspace{7\baselineskip}
Choose $\delta$ small enough that $4\delta(\psi')^2\leq\psi''/2$ throughout $[-M,0]$. By \eqref{eq:stw-concavity-square},
\[
 \frac{\mathcal C}{\lambda}
 \geq\sum_{k\in I}\frac{f_k-f_1}{\lambda(\lambda-\lambda_k)}
                         |\nabla_1g_{k\bar1}|^2
 \geq(1-2\delta)\sum_{k\in I}
                 \frac{f_k|\nabla_1g_{k\bar1}|^2}{\lambda^2}.
\]
Here $(1-\delta)/(1+\delta)\geq1-2\delta$.

\Needspace{7\baselineskip}
We need the following quantitative comparison of the two derivatives. For every $\varepsilon>0$, provided $\lambda/K$ is sufficiently large depending on the fixed function $\psi$,
\begin{equation}\label{eq:stw-derivative-cancellation}
 \sum_{k\in I}f_k\frac{|\nabla_1g_{k\bar1}|^2-|a_k|^2}{\lambda^2}
 \geq-\frac{E}{12K}-C\mathcal F
       -\varepsilon C|\psi'|\mathcal F
       -C_\varepsilon|\psi'|\sum_k f_k|v_k|^2.
\end{equation}
Here is a proof keeping the weights explicit. Let $U=\sum_{p>1,q}|v_{pq}|$. Commuting partial derivatives in $g$ gives, for $k>1$,
\[
 \nabla_1g_{k\bar1}-a_k
       =W_{k\bar1,1}-W_{1\bar1,k}+O(\lambda),\qquad
 |W_{k\bar1,1}|+|W_{1\bar1,k}|\leq C(\lambda+U).
\]
The second inequality uses the absence of $v_1$ in $Z_{k\bar1}$: differentiating in direction $1$ introduces no pure $v_{11}$. All mixed Hessians are bounded by $C\lambda$, and every other pure Hessian belongs to $U$.

\Needspace{10\baselineskip}
Use $|a+e|^2-|a|^2\geq-2|a||e|$, then \eqref{eq:stw-maximum-first}. For $k>1$, the pure Hessians in $V_k$ all belong to $U$. If $H_k=\sum_p|v_{p\bar k}|$, the resulting negative terms, divided by $\lambda^2$, are bounded by a constant times
\[
 K^{-1/2}H_k+K^{-1/2}U
       +\frac{U^2}{\lambda\sqrt K}+1
       +|\psi'||v_k|+\frac{|\psi'|}{\lambda}|v_k|U.
\]
We used $H_k\leq C\lambda$ and $\lambda\geq K$ to simplify the cross terms. The weight comparison \eqref{eq:stw-weight-comparability}, symmetry of pure derivatives and Cauchy--Schwarz give $\mathcal F U^2\leq C E$. Consequently, after multiplication by $f_k$ and summation, the first two terms are bounded by any fixed small fraction of $E/K$ plus $C\mathcal F$. The third is at most $CE/(\lambda\sqrt K)$ and is absorbed by another such fraction for large $\lambda/K$. For the last two terms use
\[
 |v_k|\leq\varepsilon+C_\varepsilon|v_k|^2,\qquad
 |v_k|U\leq |v_k|^2+U^2.
\]
The $U^2$ contribution is at most $C|\psi'|E/\lambda$, again absorbable for large $\lambda/K$. Allocating the fractions to sum to $1/12$ proves \eqref{eq:stw-derivative-cancellation}.

Insert this cancellation and the preceding inequalities into \eqref{eq:stw-key-maximum-inequality}. The remaining derivative error satisfies, by \eqref{eq:stw-maximum-first},
\[
 C\sum_k f_k|g_{1\bar1,k}|/\lambda
 \leq\frac{E}{60K}+C\mathcal F
      +\varepsilon C|\psi'|\mathcal F
      +C_\varepsilon|\psi'|\sum_k f_k|v_k|^2,
\]
after Young's inequality and adjustment of constants. Also the torsion contraction satisfies
\[
 \left|\sum_k f_kW_{k\bar k}\right|
    =\left|\sum_i t_iZ_{i\bar i}\right|
    \leq C\sum_k f_k|v_k|
    \leq\varepsilon C\mathcal F+C_\varepsilon\sum_k f_k|v_k|^2.
\]
\Needspace{9\baselineskip}
The first inequality uses the same missing gradient index in $Z_{1\bar1}$. Since $\psi'<0$ and $v_{k\bar k}=\lambda_k-\chi_{k\bar k}-W_{k\bar k}$, we conclude, decreasing the remaining coefficient of $E$ if necessary, that
\begin{equation}\label{eq:stw-first-case-closure}
 \begin{split}
 0\geq{}& f_1\left(\frac{\lambda^2}{40K}
                      -2(\psi')^2\delta^{-1}K\right)
    +\left(\frac{\psi''}{2}-C_\varepsilon|\psi'|\right)
                           \sum_k f_k|v_k|^2\\
 &-C_0\mathcal F-\varepsilon C_0|\psi'|\mathcal F
       +|\psi'|\sum_k f_k(\chi_{k\bar k}-\lambda_k).
 \end{split}
\end{equation}
In obtaining the first term we used
$|v_{1\bar1}|^2\geq\lambda^2/2-CK$, which follows from $\chi+W=O(\sqrt K)$.

For large $\lambda$, \cref{lem:stw-logdet-coercivity} bounds the last sum below by $\tau\mathcal F/2$. Choose first $\varepsilon$ with $\varepsilon C_0<\tau/4$. Next choose $D_2$ so that $\psi''/2>C_\varepsilon|\psi'|$, and then choose $D_1$ so that $|\psi'|\tau/4>C_0$ on $[-M,0]$. Finally choose $\delta$ as above. All these choices depend only on the fixed data. Equation \eqref{eq:stw-first-case-closure} then forces $\lambda^2\leq80\,\delta^{-1}\sup_{[-M,0]}|\psi'|^2K^2$, proving the desired bound in the first case.

\paragraph{The remaining eigenvalue configuration.}
Suppose now $\lambda_n<-\delta\lambda$, keeping the same constants. Since $f_n\geq\mathcal F/n$ and $|\chi+W|\leq C\sqrt K$,
\[
 \frac{E}{6K}\geq\frac{f_n}{6K}|v_{n\bar n}|^2
       \geq\frac{\delta^2\lambda^2}{12nK}\mathcal F-C\mathcal F.
\]
\Needspace{9\baselineskip}
In \eqref{eq:stw-key-maximum-inequality} discard $\mathcal C$ and the $\psi''$ term. Since $\psi$ is now fixed, $|\psi'\sum_k f_kv_{k\bar k}|\leq C\lambda\mathcal F$. Use $g_{1\bar1,k}=a_k+O(\lambda)$ and Young's inequality to bound its error by
$\frac12\sum_k f_k|a_k|^2/\lambda^2+C\mathcal F$. This leaves
\[
 0\geq-\tfrac32\sum_k f_k|a_k|^2/\lambda^2
       +\frac{\delta^2\lambda^2}{12nK}\mathcal F
       +\phi''\sum_k f_k|V_k|^2-C\lambda\mathcal F.
\]
By \eqref{eq:stw-maximum-first} and the scalar inequality
$\tfrac32|a+b|^2\leq2|a|^2+6|b|^2$,
\[
 \tfrac32\sum_k f_k|a_k|^2/\lambda^2
       \leq\phi''\sum_k f_k|V_k|^2+CK\mathcal F.
\]
As $\lambda\geq K$, cancellation and division by $\lambda\mathcal F>0$ give $\lambda/K\leq C$. This finishes both cases. The maximum property of $H$ propagates the bound to all points, and $|\lambda_i|\leq(n-1)\lambda_1$ then bounds every mixed Hessian entry, as asserted.
\end{proof}


\Needspace{10\baselineskip}
\section{A Liouville theorem for the limiting cone}
\label{sec:stw-liouville}

Rescaling at a point of large gradient removes the background and first-order terms of the Gauduchon equation and leaves a degenerate equation on $\mathbb C^n$, whose solutions are maximal for a cone of matrices. We show that its bounded Lipschitz solutions are constant. The proof argues by induction on the dimension: if the energy of a solution in some complex direction becomes small on large balls, a limit is independent of one variable; if all directional energies stay positive, they give a strictly admissible comparison function.

On $\mathbb C^n$ we use Euclidean coordinates $z=(z_1,\ldots,z_n)$, the identity matrix $I$, the Euclidean balls $B(x,r)$ and Euclidean volume $\mathrm{d}y$. In this section $\nabla$ and $\Delta$ denote the Euclidean real gradient and Laplacian, and $\partial_\xi$ denotes differentiation in a complex direction $\xi$.

\Needspace{8\baselineskip}
For $n\geq2$ and a Hermitian matrix $H$, set
\[
 P_n(H)=\frac{(\operatorname{tr}H)I-H}{n-1}.
\]
This is the map $P_\gamma$ of \cref{sec:stw-matrix-structure} for the Euclidean metric. It preserves nonnegativity and satisfies $P_n(I)=I$ and $\operatorname{tr}P_n(H)=\operatorname{tr}H$. For a real function $v$ write $H(v)=(v_{i\bar j})$ in Euclidean coordinates. A continuous function $v$ on an open set of $\mathbb C^n$ is \emph{admissible}\index{Liouville theorem!admissible functions} if $P_n(H(v))\geq0$ in the sense of matrix-valued distributions. It is \emph{maximal for this cone} if it is admissible and every smooth function $h$ touching $v$ from below at a point satisfies
\[
 P_n(H(h))\not>0
\]
at that point. Here touching means that $v-h$ has a local minimum with value zero, which need not be a strict minimum. Maximality is thus the viscosity supersolution condition for the degenerate equation $\det P_n(H(v))=0$ on the boundary of the cone, together with distributional admissibility.

Taking the trace shows that admissible functions are subharmonic. Since $P_n$ is linear, convolution preserves admissibility, and so do nonnegative linear combinations and integrals of admissible smooth functions. Translations and unitary changes of coordinates preserve both conditions, as do multiplication by a positive constant and addition of a constant.

\begin{lemma}\label{lem:cone-maximal-stability-comparison}
Local uniform limits of admissible maximal functions are admissible and maximal. If $v$ is maximal on a bounded open set $D\subset\mathbb C^n$, $h$ is smooth and admissible on $D$, both functions are continuous on $\overline D$, and $h\leq v$ on $\partial D$, then $h\leq v$ on $D$.
\end{lemma}
\begin{proof}
Let $v_j$ be admissible maximal functions converging locally uniformly to $v$. Distributional positivity is preserved under this convergence. For maximality, suppose that a smooth function $h$ touches $v$ from below at $z_0$ and $P_n(H(h)(z_0))>0$. Choose a small ball on which $v-h\geq0$. Replacing $h$ by $h-\delta|z-z_0|^2$, with $\delta>0$ sufficiently small, keeps its matrix strictly positive near $z_0$ and makes the minimum strict. Uniform convergence then produces interior minimum points of $v_j-h+\delta|z-z_0|^2$ tending to $z_0$. Adding the minimum value to the test function gives a function touching $v_j$ from below with strictly positive matrix, a contradiction.

For comparison, choose $R$ with $\overline D\subset\{|z|\leq R\}$. If $h>v$ somewhere, then for small $\delta>0$ the smooth function
\[
 h_\delta=h+\delta(|z|^2-R^2)
\]
still exceeds $v$ at some interior point and is at most $v$ on the boundary. Hence the negative minimum of $v-h_\delta$ is attained in the interior. But
$P_n(H(h_\delta))=P_n(H(h))+\delta I>0$ there, which contradicts maximality after the minimum value is added to $h_\delta$. No boundary regularity of $D$ is needed for this argument.
\end{proof}

\begin{theorem}\label{thm:stw-cone-liouville}\index{Liouville theorem!Gauduchon cone}
A bounded globally Lipschitz function on $\mathbb C^n$, $n\geq2$, that is maximal for this cone is constant.
\end{theorem}
\begin{proof}
Let $v$ be a bounded globally Lipschitz maximal function on $\mathbb C^n$. We argue by induction on $n$. For $n=2$, $P_2(H)\geq0$ is equivalent to $H\geq0$, so $v$ is plurisubharmonic, and \cref{lem:psh-liouville} shows that it is constant.

We first explain the reduction in dimension for $n\geq3$. If a function is independent of $z_n$, its complex Hessian has the block form $H\oplus0$. The first block of $P_n(H\oplus0)$ is
\[
 \frac{n-2}{n-1}P_{n-1}(H),
\]
and the last entry is $\operatorname{tr}H/(n-1)$. Nonnegativity of $P_{n-1}(H)$ implies $\operatorname{tr}H\geq0$, and positive definiteness implies $\operatorname{tr}H>0$. Thus $P_n(H\oplus0)$ is nonnegative, respectively positive definite, exactly when $P_{n-1}(H)$ is. The distributional version follows from the same linear identities. A test function in $n-1$ variables extends to $\mathbb C^n$ independently of $z_n$, so maximality also passes to $n-1$ variables. Consequently the induction hypothesis applies to bounded Lipschitz maximal functions that are independent of one complex variable.

Suppose now that $v$ is nonconstant. Adding a constant and multiplying by a positive constant, we may assume that $\inf v=0$ and $\sup v=1$. Let $L$ be the Lipschitz constant of $v$. Choose a smooth radial nonnegative probability kernel $\eta$, supported in the closed unit ball and strictly positive in its interior. For a bounded measurable function $f$ and $r>0$, write
\[
 [f]_r(x)=\int_{\mathbb C^n}f(x+y)\eta_r(y)\,\mathrm{d}y,
 \qquad \eta_r(y)=r^{-2n}\eta(y/r),\qquad v_\epsilon=[v]_\epsilon.
\]
Both $v$ and $v^2$ are subharmonic. For $v^2$, use $v_\epsilon\geq0$ and
$\Delta(v_\epsilon^2)=2v_\epsilon\Delta v_\epsilon+2|\nabla v_\epsilon|^2\geq0$, and pass to the uniform limit as $\epsilon\to0$.

For $f=v$ and for $f=v^2$, the averages $[f]_r(x)$ increase with $r$, by the monotonicity of spherical means of subharmonic functions. Their limit as $r\to\infty$ is independent of $x$, since
\[
 |[f]_r(x)-[f]_r(y)|
 \leq\|f\|_\infty\|\eta_r(\cdot-x)-\eta_r(\cdot-y)\|_{L^1}
 \leq C\frac{|x-y|}{r}.
\]
By the submean inequality this limit is at least $f(x)$ for every $x$, while it is at most $\sup f=1$. Therefore
\begin{equation}\label{eq:cone-large-averages}
 [v]_r(x)\longrightarrow1,\qquad [v^2]_r(x)\longrightarrow1.
\end{equation}

We distinguish two cases. First suppose that there exist $\rho>0$, sequences $\epsilon_k\downarrow0$ and $r_k\to\infty$, points $x_k$ and unit complex directions $\xi_k$ such that
\begin{equation}\label{eq:cone-gap-case}
 [v^2]_{r_k}(x_k)+[v]_\rho(x_k)-2v(x_k)\geq\frac43,
 \qquad
 \int_{B(x_k,r_k)}|\partial_{\xi_k}v_{\epsilon_k}|^2\,\mathrm{d}y
       \longrightarrow0.
\end{equation}
Translate $x_k$ to the origin and make $\xi_k$ the last complex direction by a unitary change of coordinates. The resulting functions have the same bounds and Lipschitz constant, so by the Arzel\`a--Ascoli theorem and \cref{lem:cone-maximal-stability-comparison} a subsequence converges locally uniformly to a maximal function $a$. The corresponding mollifications at scale $\epsilon_k$ have the same local uniform limit, and by the integral condition their derivatives in the last complex direction tend to zero in $L^2$ on every fixed ball. Hence $\partial a/\partial z_n=0$ in the sense of distributions. Since $a$ is real and continuous, it is independent of $z_n$.

The dimension reduction and the induction hypothesis show that $a$ is a constant $c\in[0,1]$. In the first inequality of \eqref{eq:cone-gap-case}, the first term is at most one, while the average at the fixed radius $\rho$ and the value at the origin both tend to $c$. It would follow that $4/3\leq1-c\leq1$, a contradiction.

\Needspace{16\baselineskip}
In the remaining case, for every fixed $\rho>0$ there is $C_\rho\geq1$ such that, whenever $0<\epsilon<C_\rho^{-1}$, $r>C_\rho$, and
\[
 [v^2]_r(x)+[v]_\rho(x)-2v(x)\geq\frac43,
\]
one has
\begin{equation}\label{eq:cone-directional-energy}
 \int_{B(x,r)}|\partial_\xi v_\epsilon|^2\,\mathrm{d}y
       \geq C_\rho^{-1}
       \quad\hbox{for every unit complex direction }\xi.
\end{equation}
Indeed, if this failed for $C_\rho=k$ for every integer $k$, a diagonal choice would give sequences as in \eqref{eq:cone-gap-case}.

After a translation we may assume that $v(0)<1/12$. By \eqref{eq:cone-large-averages}, choose $\rho$ with $[v]_\rho(0)>3/4$, and then $r>C_\rho$ with $[v^2]_r(0)>3/4$. The open set
\[
 D=\{x:2v(x)<[v^2]_r(x)+[v]_\rho(x)-4/3\}
\]
contains the origin. We claim that, for some $c>0$,
\[
 [v^2]_r-c|z|^2
\]
is smooth and admissible on $D$.

To prove the claim, differentiate the square of the smooth mollification:
\[
 H(v_\epsilon^2)
       =2v_\epsilon H(v_\epsilon)
            +2(\partial_i v_\epsilon\,\overline{\partial_jv_\epsilon})_{ij}.
\]
Denote by $A_\epsilon$ and $B_\epsilon$ the convolutions at the fixed radius $r$ of these two terms. The first satisfies $P_n(A_\epsilon)\geq0$, because $v_\epsilon\geq0$, $v_\epsilon$ is admissible, and $P_n$ is linear.

The second satisfies $B_\epsilon\geq c_0I$ on $D$ for small $\epsilon$, with a constant $c_0>0$. Here the passage from the unweighted integral in \eqref{eq:cone-directional-energy} to the convolution needs care. The bound $|\nabla v_\epsilon|\leq L$ makes the energy of a sufficiently thin boundary annulus of $B(x,r)$ at most $C_\rho^{-1}/2$, uniformly in $x,\epsilon,\xi$. On the remaining smaller closed ball, the fixed kernel $\eta_r$ has a positive minimum. For points $x\in D$, \eqref{eq:cone-directional-energy} then gives a uniform positive lower bound for the weighted integral in every direction, which is the asserted matrix inequality for $B_\epsilon$.

Choose $0<c\leq c_0$. Since $P_n$ preserves nonnegative matrices,
\[
 P_n\bigl(H([v_\epsilon^2]_r)-cI\bigr)\geq0
       \quad\hbox{on }D.
\]
As $\epsilon\downarrow0$, $v_\epsilon^2\to v^2$ uniformly, and convolution at the fixed positive radius $r$ gives convergence of all derivatives. This proves the claim.

The smooth function
\[
 h=\tfrac12\bigl([v^2]_r-c|z|^2+[v]_\rho-4/3\bigr)
\]
is therefore admissible on $D$. The set $D'=\{v<h\}$ contains the origin and lies in $D$. It is bounded, because $v\geq0$ while $h\to-\infty$ at infinity, the averages being bounded. Moreover $v=h$ on $\partial D'$. \Cref{lem:cone-maximal-stability-comparison}, applied on $D'$, gives $h\leq v$ there, which contradicts the definition of $D'$. This excludes the second case and proves the theorem.
\end{proof}

The Lipschitz bound is used both for the compactness of translates and for the annulus estimate. For $n>1$, the conclusion would not follow from boundedness and subharmonicity alone, since $\mathbb C^n=\mathbb R^{2n}$ carries nonconstant bounded subharmonic functions.


\section{The gradient bound by rescaling}
\label{sec:stw-gradient-reduction}

We now explain precisely how the second-order estimate gives the gradient bound for the Gauduchon equation. We state the implication separately, so that the role of its hypothesis remains explicit. For a family with bounded $C^2$ data, \cref{thm:stw-second-order-reconstructed} supplies that hypothesis.

Fix the smooth Hermitian forms $\gamma_0,\gamma$ of \cref{thm:stw-volume}. For a smooth solution $u$ write
\[
 \Gamma_u=\gamma_0^{n-1}+\mathcal D_\gamma u=\eta_u^{n-1}>0,
 \qquad \eta_u^n=e^{F+b}\gamma^n,\qquad \sup_Xu=0,
\]
using the operator $\mathcal D_\gamma$ defined above. Consider a family with $\|F\|_{C^0}\leq M$. The preceding $C^0$ estimate bounds $\|u\|_{C^0}$ and $|b|$ uniformly.

\Needspace{11\baselineskip}
\begin{proposition}\label{prop:stw-gradient-from-second-order}
Suppose there is a constant $K$ such that every solution in this family satisfies
\begin{equation}\label{eq:stw-second-order-required}
 \sup_X|\ddc u|_\gamma
       \leq K\bigl(1+\sup_X|\partial u|_\gamma^2\bigr).
\end{equation}
Then $\sup_X|\partial u|_\gamma$ and $\sup_X|\ddc u|_\gamma$ are uniformly bounded in the family. No existence theorem is assumed for this assertion about solutions that are already given.
\end{proposition}
\begin{proof}
We first give the matrix form needed under rescaling. Use the Hodge star of $\gamma$ to identify real $(n-1,n-1)$-forms with Hermitian matrices, normalized so that $\gamma^{n-1}$ corresponds to $I$. In a $\gamma$-unitary frame let $Q_u$ be the matrix corresponding to $\Gamma_u$. The wedge identity gives
\begin{equation}\label{eq:stw-cofactor-matrix-scaling}
 Q_u=Q_0+P_n\bigl((\ddc u)_\gamma\bigr)+Z_\gamma(\partial u),
 \qquad \det Q_u=e^{(n-1)(F+b)}.
\end{equation}
Here $Q_0>0$ is fixed and smooth, $Z_\gamma$ is real linear in the first derivatives with fixed smooth coefficients, and $(\ddc u)_\gamma$ is the Hermitian coefficient matrix of $\ddc u$. To check the principal term, diagonalize that matrix at the point: the coefficient in the complementary $(n-1)$-fold wedge is the sum of all diagonal entries except the omitted one, divided by $n-1$. Unitary covariance proves the formula for all matrices. The other term in $\mathcal D_\gamma$ contains only first derivatives of $u$, giving $Z_\gamma$. Finally, if the eigenvalues of $\eta_u$ relative to $\gamma$ are $\lambda_i$, the eigenvalues of $Q_u$ are $\prod_{j\ne i}\lambda_j$. Their product is $(\prod_i\lambda_i)^{n-1}$, proving the determinant identity.

Suppose the gradient bound fails. Choose solutions $u_j$ with
\[
 A_j=\sup_X|\partial u_j|_\gamma
       =|\partial u_j(x_j)|_\gamma\longrightarrow\infty.
\]
After extraction the points lie in one fixed coordinate neighborhood. Translate their coordinate centers to zero and make a uniformly controlled complex linear change so that $\gamma(x_j)$ is the identity. A fixed-radius coordinate ball remains available at every center. On balls whose radii tend to infinity define
\[
 v_j(z)=u_j(x_j+z/A_j),
\]
where the preceding linear changes are understood. The coefficient matrices $g_j(z)$ of $\gamma$ evaluated at these scaled points converge smoothly on every fixed ball to the Euclidean identity. In this notation
\[
 |v_j|\leq C,
 \qquad |\partial v_j|_{g_j}\leq1,
 \qquad |\partial v_j(0)|=1.
\]
The assumed estimate gives a uniform bound for all complex Hessian entries of $v_j$ on every fixed ball, since two derivatives contribute $A_j^{-2}$. Hence their Euclidean Laplacians are uniformly bounded. The standard interior estimates for the scalar Laplacian give local $W^{2,p}$ bounds for every finite $p$. Choose $p>2n$ and use Sobolev embedding and compactness to extract a diagonal subsequence converging in $C^1$ on compact subsets to a function $v$ on $\mathbb C^n$. It is bounded and globally Lipschitz, and
\[
 |\partial v(0)|=1.
\]

Divide \eqref{eq:stw-cofactor-matrix-scaling} by $A_j^2$ and evaluate at the scaled points. The background term is $O(A_j^{-2})$. The first-order term is $O(A_j^{-1})$, since one derivative of $u_j$ is $A_j$ times a derivative of $v_j$. The uniformly bounded complex Hessians and $g_j\to I$ therefore give, in fixed Euclidean coordinates on any fixed ball,
\begin{equation}\label{eq:stw-rescaled-matrix}
 A_j^{-2}Q_{u_j}
       =P_n\left(\frac{H(v_j)}{2\pi}\right)+o(1),
 \qquad
 \det(A_j^{-2}Q_{u_j})
       =A_j^{-2n}e^{(n-1)(F_j+b_j)}\longrightarrow0.
\end{equation}
For precision, take the compatible unitary frames obtained from the positive square roots $S_j=g_j^{-1/2}$. The principal matrix in these frames is $P_n(S_jH(v_j)S_j/(2\pi))$, which differs from $P_n(H(v_j)/(2\pi))$ by a uniformly vanishing error. No derivatives of these frames are taken, and the determinant identity is exact in the unitary frames. The factor $1/(2\pi)$ records the convention for $\ddc$; it does not change the cone. The positive matrices on the left imply, by passage to distributions, $P_n(H(v))\geq0$.

We check maximality. Suppose a smooth $h$ touches $v$ from below at $z_0$ and $P_n(H(h)(z_0))>0$. Subtract $\delta|z-z_0|^2$ with $\delta>0$ small. On a small ball this makes the contact strict and preserves strict positivity of the matrix. Local uniform convergence supplies interior lower contact points $z_j\to z_0$ for $v_j$ and this perturbed test function, after adding constants to the latter. The complex Hessian inequality at a minimum gives
\[
 H(v_j)(z_j)\geq H(h)(z_j)-\delta I.
\]
Because $P_n$ preserves nonnegative matrices, \eqref{eq:stw-rescaled-matrix} now implies $A_j^{-2}Q_{u_j}(z_j)\geq cI$ for a fixed $c>0$ and all sufficiently large $j$. Its determinant is at least $c^n$, contradicting the limit zero in that equation. Thus $v$ is maximal for the cone of \cref{sec:stw-liouville}.

The Liouville theorem makes $v$ constant, contradicting $|\partial v(0)|=1$. This proves the uniform gradient bound. Substitution in \eqref{eq:stw-second-order-required} gives the uniform complex Hessian bound.
\end{proof}

This is the rescaling step of \cite[Section~6]{Sze18}, with the first-order term of the Gauduchon ansatz retained explicitly. That term disappears at rate $A_j^{-1}$; it cannot be omitted in the original equation. For prescribed functions bounded in $C^2$, the potential estimate bounds $u$ and $b$, and \cref{thm:stw-second-order-reconstructed}, after the fixed conversion $v=u/(2\pi)$, proves \eqref{eq:stw-second-order-required}. Thus both the gradient and the complex Hessian are uniformly bounded. We next use these bounds to close the continuity argument.


\section{Closing the Gauduchon continuity argument}
\label{sec:stw-continuity}

We explain the remaining regularity and solvability steps using standard local elliptic theory. The relevant concavity acts on a linear transform of the complex Hessian, so its real uniformly elliptic formulation needs to retain that transform.

\begin{lemma}\label{lem:stw-higher-regularity}
For the fixed forms in \cref{thm:stw-volume}, the preceding potential, gradient and second-order estimates give uniform $C^k$ bounds for smooth normalized solutions whenever the smooth prescribed functions $F$ range over a set bounded in every $C^k$ norm.
\end{lemma}
\begin{proof}
Work first with $v=u/(2\pi)$. The second-order theorem and the rescaling proposition give uniform bounds for $v$, its gradient and its complex Hessian. The constants $b$ are bounded by the potential estimate. Hence the positive matrix
\[
 Q=Q_0+P_\gamma(\mathrm i\partial\bar\partial v)+Z(\partial v)
\]
has uniformly bounded eigenvalues above. Its determinant is uniformly bounded below by the equation. Every eigenvalue is therefore bounded below by a positive constant as well. Choose fixed $0<a<b_0$ such that $aI\leq Q\leq b_0I$ in unitary frames for every solution.

The fixed scalar operator $\Delta_\gamma v=\operatorname{tr}_\gamma(\mathrm i\partial\bar\partial v)$ is uniformly elliptic and has bounded right-hand side. Its standard interior $W^{2,p}$ estimates, applied on a finite cover, bound $v$ in $W^{2,p}$ for every finite $p$. Choosing $p>2n$ gives a uniform $C^{1,\alpha}$ bound for some $\alpha>0$. Thus the matrix $R(x)=Q_0(x)+Z_x(\partial v(x))$ is uniformly $C^\alpha$.

On a coordinate ball, let $\mathcal H_x(A)$ denote the Hermitian matrix of $\mathrm i\partial\bar\partial v$ in a fixed smooth unitary frame when $A=D^2_{\mathbb R}v$. It is a real linear map in $A$, smooth in $x$. If a real symmetric matrix $N$ is nonnegative, then $\mathcal H_x(N)\geq0$, and its trace is comparable to $\operatorname{tr}_{\mathbb R}N$ with uniform positive constants on the ball. This follows either by evaluating on a vector and its complex rotation, or by the coordinate formula for $v_{i\bar j}$.

Freeze the already obtained function $\partial v(x)$ in $R(x)$ and define on all real symmetric $A$
\[
 \widetilde\Phi(A,x)=
  \inf_{aI\leq B\leq b_0I}
    \left\{\log\det B+
       \operatorname{tr}\left[B^{-1}
          \bigl(R(x)+P_n(\mathcal H_x(A))-B\bigr)\right]\right\}.
\]
The infimum is over a fixed compact set of positive Hermitian matrices. This function is concave in $A$. If $R(x)+P_n(\mathcal H_x(A))$ lies in that set, it agrees with its logarithmic determinant: each expression is a supporting affine function of the concave logarithmic determinant, and the tangent at that very matrix attains equality.

For $N\geq0$, each affine function increases by
\[
 \operatorname{tr}\bigl(B^{-1}P_n(\mathcal H_x(N))\bigr).
\]
Since $P_n$ preserves positivity and trace, this lies between fixed positive multiples of $\operatorname{tr}_{\mathbb R}N$, uniformly in $B,x$. Taking infima preserves both inequalities. Thus $\widetilde\Phi$ is uniformly elliptic on all real symmetric matrices. Smoothness of $\mathcal H_x$ and the uniform $C^\alpha$ bound for $R$ also give
\[
 |\widetilde\Phi(A,x)-\widetilde\Phi(A,y)|
          \leq C|x-y|^\alpha(1+\|A\|).
\]
These are the hypotheses of the standard local concave, uniformly elliptic Evans--Krylov estimate with H\"older coefficients; one precise formulation is \cite[Definition~3.2 and Theorem~3.4]{TWWY15}. The value $\widetilde\Phi(0,x)$ is uniformly bounded by the compact tangent set and the bound for $R$. The given smooth function solves
\[
 \widetilde\Phi(D^2_{\mathbb R}v,x)=(n-1)(F+b)
\]
on the ball. That estimate, on smaller balls of a finite cover, yields a uniform $C^{2,\beta}$ bound for some $\beta>0$.

Differentiating the original smooth matrix equation gives linear uniformly elliptic equations for first derivatives of $v$. Their coefficients, including the first-order coefficients from $Z$, are now H\"older continuous. Local Schauder estimates give $C^{3,\beta}$ bounds, and differentiation and induction give every higher bound. The normalizations and the finite coordinate cover make the bounds global. The only regularity inputs here are the standard scalar $W^{2,p}$, concave Evans--Krylov and Schauder estimates.
\end{proof}

\begin{proof}[of \cref{thm:stw-volume} from the preceding estimates]
Let
\[
 F_0=\log\frac{\gamma_0^n}{\gamma^n},\qquad
 F_t=(1-t)F_0+tF,\quad0\leq t\leq1.
\]
At $t=0$, $u=0,b=0$ solves the equation. For the continuity argument normalize by $\int_Xu\,\gamma^n=0$; adding a constant has no effect on the equation, and at the end we impose $\sup u=0$ instead.

We verify openness including its scalar obstruction. At a solution, the derivative in $u$ of $\log\det Q_u$ is a real scalar uniformly elliptic operator $L$ with a first-order term and no zeroth-order term. Its principal positivity follows from the matrix calculation in the lemma, or directly from the positive trace pairing with $P_n$. The maximum principle on connected $X$ gives $\ker L=\mathbb R$. The standard adjoint construction used in \cref{sec:gauduchon-volume-adjoint} supplies a positive smooth probability measure $\mu$ with $L^*\mu=0$. Fredholm theory gives the range condition
\[
 Lw=f\quad\Longleftrightarrow\quad\int_X f\,\mathrm{d}\mu=0.
\]
The index is zero: the principal elliptic coefficients can be joined through positive coefficients to a fixed scalar Laplacian, whose index is zero, while the lower-order terms can be continuously removed. Its cokernel is therefore one-dimensional. For the logarithmic determinant equation the full derivative in $(u,b)$ is
\[
 (w,c)\longmapsto Lw-(n-1)c.
\]
Given any H\"older function $f$, take $c=-(n-1)^{-1}\int_Xf\,\mathrm{d}\mu$ and solve $Lw=f+(n-1)c$ with zero $\gamma^n$ mean. This gives a bounded inverse on the corresponding H\"older spaces. Positivity of $\Gamma_u$ is an open condition, so the implicit function theorem proves openness of the solution parameters.

For closedness, the prescribed path has uniform bounds in every smooth norm. The potential estimate bounds $u$ after supremum normalization and bounds $b$. The second-order theorem and the rescaling argument give the gradient and complex Hessian bounds, and the preceding lemma gives every higher bound. Subtracting the mean preserves all derivative bounds and changes the uniformly bounded potentials by uniformly bounded constants. A convergent subsequence of mean-zero solutions therefore gives a smooth solution at a limiting parameter. The determinant lower bound and matrix upper bound preserve strict positivity in the limit. Thus the nonempty open and closed set of parameters is all of $[0,1]$.

Finally consider two solutions $(u,b)$ and $(v,c)$ for the same $F$. At a maximum of $u-v$ their gradients agree and their complex Hessians are ordered. The first-order terms are equal there, and $P_n$ preserves the Hessian ordering. Consequently $Q_u\leq Q_v$, so $b\leq c$ by the equation. At a minimum the reverse inequality gives $b\geq c$. Hence $b=c$. Along the segment between $u$ and $v$, the matrices $Q$ remain positive because the ansatz is affine. Integrating the linearized logarithmic determinant along that segment produces a scalar uniformly elliptic operator with drift and no zeroth-order term annihilating $u-v$. The maximum principle makes $u-v$ constant. Supremum normalization makes the constant zero.
\end{proof}

When $\gamma_0$ is Gauduchon, the elementary Aeppli identity preceding the estimates shows that the resulting metric is Gauduchon and its $(n-1)$st power is in the specified Aeppli class. The theorem remains about the scalar ansatz: it does not give uniqueness among all metrics in that class.


\Needspace{11\baselineskip}
\section{Normalization for the Gauduchon equation}

By \cref{thm:stw-volume}, every smooth real function $F$ determines a scalar $b_\gamma(F)$, the constant $b$ of that theorem. As for the Monge--Amp\`ere equation, its variation is governed by the adjoint probability measure $\mu_u^\gamma$ of \cref{prop:gaud-volume-linearization} and by the concavity of the logarithmic determinant. To differentiate solutions, we fix a smooth positive volume form $\mathrm{d}V_0$ and use the normalization $\int_Xu\,\mathrm{d}V_0=0$ in place of $\sup_Xu=0$.

\begin{lemma}\label{lem:gaud-volume-smooth-dependence}
The solution $(u,b_\gamma(F))$, normalized by $\int_Xu\,\mathrm{d}V_0=0$, depends smoothly on the parameters of every smooth finite-dimensional family of functions $F$.
\end{lemma}
\begin{proof}
The map $\Gamma\mapsto\eta$ with $\eta^{n-1}=\Gamma$ is smooth on positive forms, since in a frame it is the inverse of the smooth map $\eta\mapsto\eta^{n-1}$, whose differential is invertible. The openness argument in the proof of \cref{thm:stw-volume}, with the $\mathrm{d}V_0$-mean in place of the $\gamma^n$-mean, gives a bounded inverse for the derivative with respect to the potential and the scalar between the relevant H\"older spaces. The implicit function theorem gives local smooth solution maps, and uniqueness shows that they agree on overlaps. Applying this for every order of H\"older regularity gives the assertion for smooth data.
\end{proof}

\Needspace{7\baselineskip}
For a smooth family of prescribed functions $F_t\in C^\infty(X,\mathbb R)$, with $t$ in an open interval of $\mathbb R$, let $u_t$ be the corresponding solutions with $\int_Xu_t\,\mathrm{d}V_0=0$, and set $Q_t=Q_{u_t}$ as in \eqref{eq:stw-cofactor-matrix-scaling}. The expression
\[
 \operatorname{tr}\bigl((Q_t^{-1}\dot Q_t)^2\bigr)
\]
is independent of the chosen frame and nonnegative, because $Q_t^{-1}\dot Q_t$ is similar to the Hermitian matrix $Q_t^{-1/2}\dot Q_t Q_t^{-1/2}$.

\Needspace{18\baselineskip}
\begin{theorem}\label{thm:gaud-volume-normalization-derivatives}
For a smooth family $F_t$ and its solutions,
\begin{align*}
 \dot b_\gamma(F_t)&=-\int_X\dot F_t\,\mathrm{d}\mu_{u_t}^\gamma,\\
 \ddot b_\gamma(F_t)&=-\int_X\ddot F_t\,\mathrm{d}\mu_{u_t}^\gamma
 -\frac1{n-1}\int_X
        \operatorname{tr}\bigl((Q_t^{-1}\dot Q_t)^2\bigr)
                          \,\mathrm{d}\mu_{u_t}^\gamma.
\end{align*}
In particular, $b_\gamma$ is concave, and along $F_t=F_0+tG$ its second derivative is strictly negative unless $G$ is constant. Moreover, for smooth real functions $F,G$ on $X$ and every real constant $a$,
\[
 b_\gamma(F+a)=b_\gamma(F)-a,
 \qquad |b_\gamma(G)-b_\gamma(F)|\leq\|G-F\|_\infty.
\]
\end{theorem}
\begin{proof}
By \eqref{eq:stw-cofactor-matrix-scaling},
\[
 \mathcal G(u)=\frac1{n-1}\log\det_\gamma Q_u.
\]
Differentiating $\mathcal G(u_t)=F_t+b_\gamma(F_t)$ once gives $\mathcal L_{u_t}\dot u_t=\dot F_t+\dot b_\gamma(F_t)$, and integration against the adjoint probability measure gives the first formula.

Because $\Gamma_u$ is affine in $u$, its second derivative along the family is $\mathcal D_\gamma\ddot u_t$. The formula for the second derivative of $\log\det$ used in the proof of \cref{thm:b-derivatives} therefore gives
\[
 \mathcal L_{u_t}\ddot u_t
 -\frac1{n-1}\operatorname{tr}\bigl((Q_t^{-1}\dot Q_t)^2\bigr)
 =\ddot F_t+\ddot b_\gamma(F_t).
\]
Integrating against $\mu_{u_t}^\gamma$ at this fixed parameter gives the second formula; no derivative of the measure is needed.

For affine $F_t$, we have $\ddot F_t=0$, and positivity of the density of $\mu_{u_t}^\gamma$ proves concavity. If the second derivative vanishes, then $\dot Q_t=0$, so $\dot\Gamma_{u_t}=0$ and $\mathcal L_{u_t}\dot u_t=0$. By \cref{prop:gaud-volume-linearization}, $\dot u_t$ is constant, and the first differentiated equation forces $G$ to be constant. Conversely, adding a constant to $F$ changes only $b$, by the opposite constant. Finally, integrate the first derivative along the segment from $F$ to $G$: since $\mu_{u_t}^\gamma$ is a probability measure, the derivative lies between $-\sup_X(G-F)$ and $-\inf_X(G-F)$, which proves the Lipschitz bound.
\end{proof}

\Needspace{6\baselineskip}
The Lipschitz bound extends $b_\gamma$ uniquely to continuous functions $F$. Constructing corresponding weak solutions would require, in addition, a volume measure for the $(n-1)$st roots of singular forms in the Aeppli ansatz, compactness of their potentials, and stability of the adjoint equation.

\Needspace{13\baselineskip}
On a surface with Gauduchon background form $\gamma_0$, conservation of the total mass makes the scalar and the adjoint probability measure explicit.

\Needspace{9\baselineskip}
\begin{example}\label{ex:gaud-volume-surface}
If $n=2$ and $\gamma_0$ is Gauduchon, then $\Gamma_u=\gamma_0+\ddc u=\eta_u$, and $\eta_u$ is also Gauduchon. The mass identity of \cref{thm:mass-invariance-characterization} gives $\int_X\eta_u^2=\int_X\gamma_0^2$. Therefore
\[
 b_\gamma(F)=\log\int_X\gamma_0^2-
                   \log\int_X\mathrm{e}^F\gamma^2,
 \qquad
 \mathrm{d}\mu_u^\gamma=\frac{\mathrm{e}^F\gamma^2}{\int_X\mathrm{e}^F\gamma^2}.
\]
The adjoint equation of \cref{prop:gaud-volume-form-adjoint} gives the same measure, because $\ddc\eta_u=0$ and hence $\rho_u$ is constant. Along $F+tG$, the second derivative of $b_\gamma$ is minus the variance of $G$ with respect to this probability measure.
\end{example}

